\chapter{Image déterminantale et composition vers Birch--Swinnerton--Dyer}

\section{Unité généalogique et produit fibré}

Soit $E/\mathbb Q$ une courbe elliptique et soit
$\mathfrak G_{\rm det}(E)$ l'image déterminantale autonome complète. On
note $\mathcal T_m^{\rm surv}$ sa catégorie d'histoires
survivantes. Le propriétaire déterminantal précédent construit à partir d'une
unité commune la chaîne de lignes
\eqref{eq:amp-bsd-linea-determinante}--\eqref{eq:amp-bsd-cadena-lineas}, le
comparateur Kato--FERS
\eqref{eq:amp-bsd-phi-kato-fers}--\eqref{eq:amp-bsd-inversa-kato-fers}, le
canal Poitou--Tate
\eqref{eq:amp-bsd-comparador-pt-q}--\eqref{eq:amp-bsd-comparador-pt} et le
recollement local--global
\eqref{eq:amp-bsd-triangulo-local}--\eqref{eq:amp-bsd-mapa-localizacion}.
Le tuple de \eqref{eq:pdf2-g-det-explicito} et les identités de
\eqref{eq:pdf2-g-det-identidades-dato} sont la publication conjointe de cette
construction en amont : ils ne sont pas acceptés comme axiomes terminaux et ne sont pas choisis à
partir du rang ou du terme principal. Ce chapitre conserve et compose
ces flèches jusqu'à la conclusion de Birch--Swinnerton--Dyer.

L'objet des coefficients n'est pas seulement la fibre $3$-primaire. On fixe d'abord
le complexe libre intégral basé
\begin{equation}\label{eq:pdf2-bsd-gmz}
 G_{m,\mathbb Z}:=N_*\bigl(N\mathcal T_m^{\rm surv};\mathbb Z\bigr),
\end{equation}
et, pour chaque nombre premier $\ell$ et chaque $n\geq1$, on pose
\begin{equation}\label{eq:pdf2-bsd-coeficientes-todos-ell}
 A_{\ell,n}:=E[\ell^n](\overline{\mathbb Q}),\qquad
 T_\ell E:=\varprojlim_nA_{\ell,n},\qquad
 V_\ell E:=T_\ell E\otimes_{\mathbb Z_\ell}\mathbb Q_\ell,
 \qquad
 G_{m,\ell,n}:=G_{m,\mathbb Z}\otimes_{\mathbb Z}^{\mathbf L}A_{\ell,n}.
\end{equation}
Les transitions ne restent pas implicites : ce sont
\begin{equation}\label{eq:pdf2-bsd-transiciones-ell-n}
 q_{\ell,n}:A_{\ell,n+1}\twoheadrightarrow A_{\ell,n},
 \quad x\longmapsto\ell x,
 \qquad
 Q_{\ell,n}:=1\otimes q_{\ell,n}:G_{m,\ell,n+1}\to G_{m,\ell,n},
\end{equation}
La transition de coefficients dans le complexe de Manin, avec
$Q_{\ell,n}$, induit
$\widetilde Q_{\ell,n}:P_{E,\ell,n+1}^{\rm gen}\to
P_{E,\ell,n}^{\rm gen}$. Ces applications satisfont
$Q_{\ell,n}d=dQ_{\ell,n}$,
$(Q_{\ell,n}\otimes Q_{\ell,n})\Delta_{\rm AW}
=\Delta_{\rm AW}Q_{\ell,n}$ et
$\widetilde Q_{\ell,n}\widehat\Pi_{0,\ell,n+1}
=\widehat\Pi_{0,\ell,n}Q_{\ell,n}$. Ainsi, la limite qui définit $T_\ell E$
est prise dans la même famille de complexes, non seulement dans les groupes de
coefficients.
L'ancienne notation $A_n$ désigne uniquement
$A_{3,n}$ ; elle ne gouverne pas la publication arithmétique. Les matrices de bord,
d'Alexander--Whitney, d'orientation et de retenue sont écrites sur $\mathbb Z$ dans
la base des histoires. Ainsi, pour
$R\in\{\mathbb Z_\ell,\mathbb Z/\ell^n,\mathbb Q_\ell\}$,
\begin{equation}\label{eq:pdf2-bsd-base-change-integral}
 G_{m,R}=G_{m,\mathbb Z}\otimes_{\mathbb Z}^{\mathbf L}R,
 \qquad
 d_{m,R}=d_{m,\mathbb Z}\otimes1,
 \qquad
 \Delta_{{\rm AW},R}=\Delta_{{\rm AW},\mathbb Z}\otimes1.
\end{equation}
Ainsi, le passage de $3$ à un autre nombre premier n'est pas une extension inexistante
$\mathbb Z_3\to\mathbb Z_\ell$ : toutes les fibres s'obtiennent par changement de
base du même propriétaire intégral.

Pour abréger, les formules suivantes sont écrites sur un coefficient
$A_{\ell,n}$ fixe et sont compatibles avec $n$, avec $\ell$ et avec la base
intégrale. Le complexe normalisé
\begin{equation}\label{eq:pdf2-bsd-gmn}
 G_{m,\ell,n}:=N_*\bigl(N\mathcal T_m^{\rm surv};A_{\ell,n}\bigr)
\end{equation}
conserve la diagonale d'Alexander--Whitney
\begin{equation}\label{eq:pdf2-bsd-aw}
 \Delta_{\rm AW}[x_0\to\cdots\to x_q]
 =\sum_{j=0}^q[x_0\to\cdots\to x_j]\otimes
                [x_j\to\cdots\to x_q]
\end{equation}
et la counité
$\varepsilon_{\rm AW}[x]=1$ sur les sommets, zéro en degrés positifs. L'histoire
se couple au complexe de Manin au moyen de
\begin{equation}\label{eq:pdf2-bsd-pullback}
 P_{E,\ell,n}^{\rm gen}:=
 P_{E,\ell,n}^{\rm Man}\times_{A_{\ell,n}[0]}G_{m,\ell,n}.
\end{equation}
Si $s_{\ell,n}:A_{\ell,n}[0]\to P_{E,\ell,n}^{\rm Man}$ est la section
basée, alors
\[
 \widehat\Pi_{0,\ell,n}(c)=
 (s_{\ell,n}\varepsilon_{\rm AW}(c),c)
\]
conserve simultanément la coordonnée modulaire et l'histoire. Cette partie de la
construction est matérielle et n'utilise pas le terme principal de $L(E,s)$. Les
carrés
\begin{equation}\label{eq:pdf2-bsd-pullback-base-change}
 P_{E,\mathbb Z}^{\rm gen}\otimes_{\mathbb Z}^{\mathbf L}
      \mathbb Z/\ell^n
 \xrightarrow{\ \sim\ }P_{E,\ell,n}^{\rm gen},
 \qquad
 \widehat\Pi_{0,\mathbb Z}\otimes1=\widehat\Pi_{0,\ell,n}
\end{equation}
prouvent qu'unité, pullback et counité sont les mêmes pour tous les
coefficients.

Les deux publications intégrales et leurs changements de base
\begin{equation}\label{eq:pdf2-bsd-copub}
 P_{E,R}=(P_{E,R}^K,P_{E,R}^{\rm PT}),
 \qquad P_{E,R}^K(1_E)=z_{K,R},\quad
 P_{E,R}^{\rm PT}(1_E)=z_{{\rm PT},R}
\end{equation}
doivent provenir de la même unité $1_E$. La dualité réduite est typée par
\[
 D_3(M):=\operatorname{RHom}_\Lambda(M^\iota,\Lambda)[-3],
 \qquad
 X^{\rm orth}:C_E^{\rm red}\xrightarrow{\sim}D_3(C_E^{\rm red}).
\]
Avec $L_{{\rm root},\ell,n}=A_{\ell,n}z_{{\rm PT},\ell,n}$ et le projecteur de complexes
$p_{\rm root}$, la normalisation basée donne
\begin{equation}\label{eq:pdf2-bsd-raiz}
 p_{\rm root}z_{K,\ell,n}=z_{{\rm PT},\ell,n}.
\end{equation}

\section{Remplisseur local et décomposition globale sans effacer le réseau fini}

La décomposition s'écrit d'abord sur le propriétaire intégral. Pour
chaque étiquette finie--singulière $\lambda=(\ell,m,n,\ldots)$, il existe un
coefficient d'incidence $a_\lambda=\ell c_\lambda\in\mathbb Z$ fixé par
$\operatorname{Rel}_{\rm det}$ ; dans la fibre correspondante,
$A_{\ell,n}$ agit par réduction de ce même entier. Le bloc libre
intégral avec $g$ primitif est
\begin{equation}\label{eq:pdf2-bsd-bloque-normal}
 \partial f=a_\lambda e+b,\qquad
 \partial e=g,\qquad
 \partial b=-a_\lambda g.
\end{equation}
Le tensoriser avec la fibre indiquée par $\lambda$ restitue littéralement
$a_\lambda=\ell c_\lambda$. En conséquence, si
$z_K=ar+ue+vb$ est un cycle et $z_{\rm PT}=r$, alors
$\partial z_K=(u-a_\lambda v)g=0$ implique $u=a_\lambda v$. L'égalité racine
\eqref{eq:pdf2-bsd-raiz} donne $a=1$ et donc,
\begin{equation}\label{eq:pdf2-bsd-filler-local}
 h_*:=-vf,\qquad \partial h_*=z_{\rm PT}-z_K.
\end{equation}

La multisection $\Sigma_{\rm det}^{\rm mult}$ étiquette chaque générateur par
\begin{equation}\label{eq:pdf2-bsd-etiqueta-completa}
 \lambda=(\ell,m,n,\gamma,\operatorname{or},\operatorname{Carr},
 \operatorname{Mem},\partial,\tau^{\rm term},\operatorname{loc}).
\end{equation}
L'ordre basé sépare trois classes disjointes et exhaustives : l'unité racine,
les blocs finis--singuliers avec terminal $h_*$ et les cellules du réseau
intégral local--global. Comme $\partial_{\rm det}$ conserve toutes les
étiquettes sauf la descente prescrite par la retenue, ces classes définissent
des projecteurs de complexes $p_{\rm root}$, $p_{\rm FS}$ et $p_{\rm lat}$ avec
\begin{equation}\label{eq:pdf2-bsd-tres-proyectores}
 p_{\rm root}+p_{\rm FS}+p_{\rm lat}=I,
 \qquad p_ip_j=\delta_{ij}p_i.
\end{equation}

\begin{theorem}[Décomposition globale basée]
\label{thm:pdf2-bsd-descomposicion-global}
Il existe un isomorphisme de complexes compatible avec les places, l'orientation,
la corestriction et la dualité
\begin{equation}\label{eq:pdf2-bsd-descomposicion-global}
 \Theta_E:\mathcal C_E^{\rm loc,red}
 \xrightarrow{\ \sim\ }
 \mathcal C_E^{\rm FS,ctr}\oplus\mathcal L_E,
\end{equation}
où
\begin{equation}\label{eq:pdf2-bsd-sumando-fs-y-lattice}
 \begin{aligned}
 \mathcal C_E^{\rm FS,ctr}
  &:=\bigoplus_{\lambda\in\Lambda_E^{\rm FS}}
   [\mathbb Zf_\lambda\to
    \mathbb Ze_\lambda\oplus\mathbb Zb_\lambda
                    \to\mathbb Zg_\lambda],\\
 \mathcal L_E&:=[M_E\xrightarrow{D_E}N_E].
 \end{aligned}
\end{equation}
Chaque bloc du premier terme possède la différentielle intégrale
\eqref{eq:pdf2-bsd-bloque-normal} ; sa fibre $(\ell,n)$ s'obtient par
$-\otimes_{\mathbb Z}^{\mathbf L}A_{\ell,n}$. Le complexe $\mathcal L_E$
conserve, sans
les contracter, les diviseurs de Smith et les composantes
Tate--Shafarevich--Tamagawa.
\end{theorem}

\begin{proof}
Les générateurs sont ordonnés par l'étiquette
\eqref{eq:pdf2-bsd-etiqueta-completa}. L'unité commune occupe le premier
bloc et est éliminée par $q_{\rm root}$. Pour chaque terminal
$(z_{\rm PT},h_*)$, les relations de $\operatorname{Rel}_{\rm det}$
identifient exactement quatre générateurs
$(f_\lambda,e_\lambda,b_\lambda,g_\lambda)$ ; la matrice d'incidence est
\eqref{eq:pdf2-bsd-bloque-normal}, avec $g_\lambda$ primitif car son
étiquette d'arrivée apparaît une seule fois dans la base orientée. Les cellules
restantes portent une étiquette locale--globale et forment les deux réseaux libres
$M_E,N_E$. Il n'y a pas de quatrième classe : la multisection est complète et
$\operatorname{Surv}_{\rm det}$ exige de chaque histoire soit un terminal
FS, soit une coordonnée de réseau. Le changement de bases est unitriangulaire à
diagonale un, car un bord n'introduit que des étiquettes antérieures. Cela
construit $\Theta_E$ et prouve simultanément son exhaustivité et sa
compatibilité avec les bases. La dualité $X^{\rm orth}$ échange les
deux réseaux et conserve chaque bloc FS, de sorte qu'elle entrelace aussi
$\Theta_E$.
\end{proof}

Sur un bloc normal, définissons
\begin{equation}\label{eq:pdf2-bsd-hfs-generadores}
 H_\lambda(g_\lambda)=e_\lambda,
 \qquad H_\lambda(b_\lambda)=f_\lambda,
 \qquad H_\lambda(e_\lambda)=H_\lambda(f_\lambda)=0.
\end{equation}
Un calcul sur les quatre générateurs donne
$\partial H_\lambda+H_\lambda\partial=I$. Par conséquent,
\begin{equation}\label{eq:pdf2-bsd-hfs-global}
 H_{\rm FS}:=\Theta_E^{-1}
 \left(\bigoplus_{\lambda\in\Lambda_E^{\rm FS}}H_\lambda\oplus0\right)
 \Theta_E,
 \qquad
 \partial H_{\rm FS}+H_{\rm FS}\partial=p_{\rm FS}.
\end{equation}
La coordonnée terminale de $\mathfrak G_{\rm det}(E)$ assure
$q_{\rm root}z_K\in\operatorname{im}p_{\rm FS}$ ; alors
$h_*=-H_{\rm FS}(q_{\rm root}z_K)$ reproduit
\eqref{eq:pdf2-bsd-filler-local}. Il est essentiel que le membre droit de
\eqref{eq:pdf2-bsd-hfs-global} soit $p_{\rm FS}$, non $q_{\rm root}$ :
contracter le réseau $\mathcal L_E$ effacerait les facteurs finis que BSD
doit publier.

\section{Propriétaire intégral et changement de base vers tous les nombres premiers}

La décomposition précédente s'exécute d'abord dans la base intégrale des
histoires, non dans une fibre primaire. Pour
\[
 R\in\{\mathbb Z,\mathbb Z_\ell,\mathbb Z/\ell^n,
          \mathbb Q_\ell\}
\]
sont fixés
\begin{equation}\label{eq:pdf2-bsd-complejos-base-change}
 \mathcal C_{E,R}^{\rm loc,red}:=
 \mathcal C_{E,\mathbb Z}^{\rm loc,red}
       \otimes_{\mathbb Z}^{\mathbf L}R,
 \qquad
 \mathcal L_{E,R}:=
 [M_E\otimes R\xrightarrow{D_E\otimes1}N_E\otimes R].
\end{equation}
L'ordre unitriangulaire des étiquettes utilisé dans
\cref{thm:pdf2-bsd-descomposicion-global} a des coefficients entiers. En
conséquence,
\begin{equation}\label{eq:pdf2-bsd-theta-base-change}
 \Theta_{E,R}=\Theta_{E,\mathbb Z}\otimes1,
 \qquad
 H_{{\rm FS},R}=H_{{\rm FS},\mathbb Z}\otimes1,
 \qquad
 \partial H_{{\rm FS},R}+H_{{\rm FS},R}\partial=p_{{\rm FS},R}.
\end{equation}
En particulier, la fibre $3$-primaire est l'une de ces spécialisations et ne
peut pas masquer les fibres $\ell\ne3$.

Avant de former un conoyau fini, on construit l'inverse rationnel. L'ordre
des terminaux de $\Sigma_{\rm det}^{\rm mult}$ donne des bases
$(m_1,\ldots,m_t)$ de $M_E$ et $(n_1,\ldots,n_t)$ de $N_E$ et une bijection
basée $\sigma_E$ entre les deux listes. Si
\[
 d_j:=\bigl\langle n_{\sigma_E(j)}^\vee,D_Em_j\bigr\rangle,
\]
la relation terminale et la dualité $X^{\rm orth}$ donnent $d_j\ne0$ avant de
prendre les dimensions. Définissons
\begin{equation}\label{eq:pdf2-bsd-d0-nilpotente}
 \begin{aligned}
  D_{E,0}:M_E\otimes\mathbb Q&\longrightarrow N_E\otimes\mathbb Q,
  &D_{E,0}m_j&=d_jn_{\sigma_E(j)},\\
  N_E^{\rm tri}&:=D_{E,0}^{-1}(D_E\otimes\mathbb Q)-I.
 \end{aligned}
\end{equation}
Un bord ne contient que son terminal et des étiquettes strictement antérieures ;
$N_E^{\rm tri}$ est donc strictement triangulaire dans la liste ordonnée et
$(N_E^{\rm tri})^t=0$. L'opérateur
\begin{equation}\label{eq:pdf2-bsd-j-e-inverso-racional}
 \boxed{
 J_E:=\left(\sum_{r=0}^{t-1}(-N_E^{\rm tri})^r\right)D_{E,0}^{-1}:
 N_E\otimes\mathbb Q\longrightarrow M_E\otimes\mathbb Q }
\end{equation}
se calcule donc avec un nombre fini d'opérations sur la matrice
d'incidence. Comme
$D_E\otimes\mathbb Q=D_{E,0}(I+N_E^{\rm tri})$, la somme géométrique finie
prouve les deux identités
\begin{equation}\label{eq:pdf2-bsd-j-e-dos-lados}
 J_E(D_E\otimes\mathbb Q)=I_{M_E\otimes\mathbb Q},
 \qquad
 (D_E\otimes\mathbb Q)J_E=I_{N_E\otimes\mathbb Q}.
\end{equation}
C'est la contraction rationnelle du réseau résiduel ; elle n'utilise pas le rang,
$\operatorname{Sha}(E)$, les ordres des groupes ni le terme principal.

La matrice intégrale $D_E$ conserve tous ses diviseurs élémentaires. Pour chaque
nombre premier, on définit, sans le déduire de la fibre $3$-primaire,
\begin{equation}\label{eq:pdf2-bsd-f-ell-def}
 D_{E,\ell}:=D_E\otimes\mathbb Z_\ell,
 \qquad
 F_{E,\ell}:=\operatorname{coker}D_{E,\ell}.
\end{equation}
Les deux identités de \eqref{eq:pdf2-bsd-j-e-dos-lados} impliquent que
$F_E:=\operatorname{coker}D_E$ est fini. Sa
décomposition primaire est donc finie et exacte :
\begin{equation}\label{eq:pdf2-bsd-f-descomposicion-primaria}
 F_E=\operatorname{coker}D_E
 \xrightarrow{\ \sim\ }
 \bigoplus_{\ell}F_{E,\ell},
 \qquad
 [n]\longmapsto([n\otimes1]_\ell)_\ell.
\end{equation}
L'inverse s'obtient générateur par générateur par le théorème chinois des restes
appliqué à chaque diviseur de Smith. C'est l'arête qui permet de publier
simultanément $\operatorname{Sha}(E)$ et tous les facteurs de Tamagawa.

\section{Bocksteins de toutes les pages}

Soit
\begin{equation}\label{eq:pdf2-bsd-s-universal}
 S_{E,0}(T):V_0[[T]]\longrightarrow W_0[[T]]
\end{equation}
la différentielle basée rationnelle construite par les matrices intégrales du
lecteur déterminantal. Pour chaque nombre premier $\ell$ et pour la carte complexe, on
obtient par changement de base
\begin{equation}\label{eq:pdf2-bsd-s-todos-ell}
 S_{E,\ell}(T):=S_{E,0}(T)\otimes_{\mathbb Q}\mathbb Q_\ell,
 \qquad
 S_{E,\mathbb C}(T):=S_{E,0}(T)\otimes_{\mathbb Q}\mathbb C.
\end{equation}
Elle est inversible sur $\mathbb Q((T))$ et donc sur chaque
$\mathbb Q_\ell((T))$ et $\mathbb C((T))$. Nous écrirons $S_E(T)$ lorsque la
base est sans importance. Comme $K[[T]]$ est un anneau de valuation discrète pour
$K=\mathbb Q,\mathbb Q_\ell$ ou $\mathbb C$, sa forme de Smith est
\begin{equation}\label{eq:pdf2-bsd-smith-t-adico}
 U(T)S_E(T)V(T)=\operatorname{diag}
 (T^{a_1},\ldots,T^{a_s},1,\ldots,1),
 \qquad a_i\ge1.
\end{equation}
Les exposants $a_i$ sont les mêmes dans toutes les fibres : ce sont les valuations
en $T$ des mineurs non nuls de la matrice rationnelle
\eqref{eq:pdf2-bsd-s-universal} ; remplacer $\mathbb Q$ par
$\mathbb Q_\ell$ ou $\mathbb C$ ne change pas ces valuations.
On obtient
\begin{equation}\label{eq:pdf2-bsd-orden-smith}
 d:=\operatorname{ord}_T\det S_E(T)
 =\sum_{i=1}^sa_i
 =\sum_{q\ge1}\dim_K V_q.
\end{equation}
La filtration produit $A_q=V_q/V_{q+1}$,
$B_q=\operatorname{im}\beta_q$ et des isomorphismes
\begin{equation}\label{eq:pdf2-bsd-bockstein-reducido}
 \bar\beta_q:A_q\xrightarrow{\sim}B_q,
 \qquad
 j_q:B_q\xrightarrow{\sim}A_q^\vee\otimes(I^q/I^{q+1}).
\end{equation}
Avec la page régulière $\bar S_E(0)$ et les signes de Koszul, toutes les
pages se recollent dans
\begin{equation}\label{eq:pdf2-bsd-rboc-lt}
 R_{\rm Boc}^{\rm der}(S_E)
 =\tau_{\rm reg}\otimes\bigotimes_{q\ge1}\det(\bar\beta_q)
 =\left[T^{-d}\det S_E(T)\right]_{T=0}.
\end{equation}
Aucune page n'est remplacée par l'ordre total $d$.

\section{Table complète Kato--FERS avant le déterminant}

Pour une histoire non dégénérée
$b_\lambda=[x_0\to\cdots\to x_i]$, soient
$b_{\lambda,\le j}=[x_0\to\cdots\to x_j]$ et
$b_{\lambda,\ge j}=[x_j\to\cdots\to x_i]$. Le poids complet du lecteur
de retenue est
\begin{equation}\label{eq:pdf2-bsd-peso-t}
 \nu_T(\gamma):=\operatorname{dep}(\gamma)
                 +\operatorname{Carr}(\gamma),
\end{equation}
de sorte que $\nu_T([x_i])=0$ et tout chemin normalisé de longueur positive
a $\nu_T\ge1$. La base $\mathcal B_i^K$ contient une ligne pour chaque
étiquette complète $\lambda$ ; la base $\mathcal B_i^{\rm Iw}$ contient le
générateur $\upsilon_i(b_\lambda)$ avec les mêmes faces, feuillet, orientation,
mémoire et terminal. Ainsi, $\upsilon_i$ est une bijection écrite par la
même liste d'étiquettes, non choisie au moyen d'un dénombrement ultérieur.

L'application de chaînes est
\begin{equation}\label{eq:pdf2-bsd-phi-kato}
 \Phi_{K/{\rm FERS}}:=
 \mu_{\rm Iw}\circ(P_E^K\widehat\otimes\operatorname{car}_T)
 \circ\Delta_{\rm AW}:
 \mathcal C_E^K\longrightarrow\mathcal C_E^{\rm Iw}.
\end{equation}
Sa table exhaustive, paramétrée par tous les $i$, toutes les étiquettes et
toutes les coupures, est
\begin{equation}\label{eq:pdf2-bsd-tabla-kato-fers}
\begin{array}{c|c|c|c}
 \text{ligne}&\text{coupure}&\text{image}&\text{poids }T\\ \hline
 b_\lambda&j=i&\upsilon_i(b_\lambda)&0\\
 b_\lambda&0\le j<i&
 \epsilon_{\lambda,j}\,
 \mu_{\rm Iw}\!\left(P_E^K(b_{\lambda,\le j}),
       \operatorname{car}_T(b_{\lambda,\ge j})\right)&
 \nu_T(b_{\lambda,\ge j})\ge1\\
 \text{simplexe dégénéré}&\text{tout }j&0&--
\end{array}
\end{equation}
où $\epsilon_{\lambda,j}$ est le signe déjà fixé par
$\operatorname{or}_{\rm det}$. De manière équivalente,
\begin{equation}\label{eq:pdf2-bsd-expansion-kato-fers}
 \Phi_{K/{\rm FERS}}^i(b_\lambda)
 =\upsilon_i(b_\lambda)+
  \sum_{j=0}^{i-1}T^{\nu_T(b_{\lambda,\ge j})}
  \epsilon_{\lambda,j}\,
  \mu_{\rm Iw}\!\left(P_E^K(b_{\lambda,\le j}),
          \widehat{\operatorname{car}}_T(b_{\lambda,\ge j})\right).
\end{equation}

\begin{theorem}[Équivalence basée Kato--FERS]
\label{thm:pdf2-bsd-equivalencia-kato-fers}
L'application \eqref{eq:pdf2-bsd-phi-kato} est une équivalence filtrée de
complexes parfaits, préserve l'unité et, dans toutes les bases complètes,
\begin{equation}\label{eq:pdf2-bsd-i-plus-tb}
 [\Phi_{K/{\rm FERS}}^i](T)=I+TB_i(T),
 \qquad
 [\Phi_{K/{\rm FERS}}^i]^{-1}
 =\sum_{r\ge0}(-TB_i(T))^r.
\end{equation}
En particulier,
\begin{equation}\label{eq:pdf2-bsd-rho-uno}
 \rho_E(T):=
 \prod_i\det(I+TB_i(T))^{(-1)^i},
 \qquad \rho_E(0)=1,
 \qquad
 \rho_{E,\ell}:=\rho_E\otimes_{\mathbb Q}\mathbb Q_\ell.
\end{equation}
\end{theorem}

\begin{proof}
L'identité simpliciale d'Alexander--Whitney, la compatibilité de $P_E^K$
avec les faces et l'additivité de $\nu_T$ sous concaténation prouvent que
$\Phi_{K/{\rm FERS}}$ commute avec les différentielles. Dans chaque ligne de
\eqref{eq:pdf2-bsd-tabla-kato-fers}, l'unique terme de poids nul est la
coupure $j=i$ ; tous les autres appartiennent à $T\Lambda$. Comme
$\upsilon_i$ apparie la liste complète des histoires avec elle-même, la
réduction modulo $T$ est l'identité sur chaque générateur. Cela donne
\eqref{eq:pdf2-bsd-i-plus-tb} ; la complétude de $K[[T]]$ fait converger
l'inverse. Pour $i=0$, la table ne contient que l'unité, de sorte que la racine
est normalisée avant l'évaluation d'un terme principal.
\end{proof}

L'équivalence conserve séparément le ledger
\begin{equation}\label{eq:pdf2-bsd-ledger-diez}
 \mathcal J_{\rm det}=\{\deg,\operatorname{Frob},\operatorname{Ner},
 K,\operatorname{Boc},\operatorname{Sha},\operatorname{Tam},
 \operatorname{tors}^{-2},\Omega,\operatorname{exc}\}.
\end{equation}
Dans ces dix étiquettes, on lit respectivement : le poids total ; la norme
d'Euler ; le changement basé Manin--Néron ; la relation finie--singulière ; toutes les
pages $j_q\bar\beta_q$ ; le cône de Kummer ; les quotients de composantes ;
le réseau torsionnel ; l'intégrale de Néron ; et l'orientation exceptionnelle. La
table ne permet pas d'annulations entre étiquettes.

\section{Application Poitou--Tate de Mordell--Weil vers le drapeau complet}

On conserve d'abord le réseau de Mordell--Weil et on ne rationalise qu'ensuite :
\begin{equation}\label{eq:pdf2-bsd-mw-libre}
 M_{E,\mathbb Z}^{\rm MW}:=E(\mathbb Q)/E(\mathbb Q)_{\rm tors},
 \qquad
 M_E^{\rm MW}:=M_{E,\mathbb Z}^{\rm MW}\otimes_{\mathbb Z}K,
 \qquad K=\mathbb Q.
\end{equation}
$\operatorname{Surv}_{{\rm det},\mathbb Z}^{\rm MW}$ désigne le sous-module
libre engendré par les lignes terminales intégrales du terme
Mordell--Weil.
Le relèvement de Kummer avec histoire complète est
\begin{equation}\label{eq:pdf2-bsd-kummer-hmt}
 \begin{aligned}
 \kappa_{E,\mathbb Z}^{\rm HMT}:M_{E,\mathbb Z}^{\rm MW}
   &\longrightarrow
   Z^1(\mathcal C_{E,\mathbb Z}^{\rm red})
   \cap\operatorname{Surv}_{{\rm det},\mathbb Z}^{\rm MW},\\
 \kappa_E^{\rm HMT}&:=
   \kappa_{E,\mathbb Z}^{\rm HMT}\otimes_{\mathbb Z}\mathbb Q:
   M_E^{\rm MW}\longrightarrow Z^1(\mathcal C_E^{\rm red}).
 \end{aligned}
\end{equation}
et $\operatorname{pg}_q$ publie la coordonnée de la page $V_q$.
L'extension $\operatorname{Ext}_{\Gamma,{\rm det}}$, la multisection et la
survie construisent, pour chaque $q$, un recollement basé
\begin{equation}\label{eq:pdf2-bsd-glue-q}
 \operatorname{gl}_q:V_q\longrightarrow
 \operatorname{Surv}_{\rm det}^{\rm MW}
\end{equation}
qui conserve feuillet, retenue, mémoire et terminal. Ses équations
d'unité--counité sont
\begin{equation}\label{eq:pdf2-bsd-unidad-counidad-paginas}
 \operatorname{pg}_rP_E^{\rm PT}\operatorname{gl}_q
 =\delta_{rq}I_{V_q},
 \qquad
 \sum_{q\ge1}\operatorname{gl}_q
     \operatorname{pg}_qP_E^{\rm PT}=I
 \quad\text{dans le terme MW survivant}.
\end{equation}
La première égalité se vérifie étiquette par étiquette : le recollement de niveau
$q$ a la retenue $q$, zéro dans les autres pages et le terminal de cette même
histoire. La deuxième est l'exhaustivité de
$\Sigma_{\rm det}^{\rm mult}$ ; la somme est finie car la tour de Bockstein
se termine.

Pour ce sous-module, posons
\begin{equation}\label{eq:pdf2-bsd-redes-paginas-saturadas}
 V_{q,\mathbb Z}:=\operatorname{pg}_qP_E^{\rm PT}
   (\operatorname{Surv}_{{\rm det},\mathbb Z}^{\rm MW}),
 \qquad
 A_{q,\mathbb Z}:=V_{q,\mathbb Z}/V_{q+1,\mathbb Z}.
\end{equation}
Dans les bases terminales, $\operatorname{pg}_q$ et
$\operatorname{gl}_q$ ont des entrées $0,\pm1$ et leurs deux compositions sont
celles de \eqref{eq:pdf2-bsd-unidad-counidad-paginas}. Par conséquent,
$V_{q+1,\mathbb Z}$ est un facteur direct basé de $V_{q,\mathbb Z}$,
$A_{q,\mathbb Z}$ est libre et saturé, et
$V_q=V_{q,\mathbb Z}\otimes\mathbb Q$,
$A_q=A_{q,\mathbb Z}\otimes\mathbb Q$.
Soit $s_q:A_{q,\mathbb Z}\to V_{q,\mathbb Z}$ la section de la base
terminale. Alors
$V_{k,\mathbb Z}=\bigoplus_{q\geq k}s_q(A_{q,\mathbb Z})$. Le dédoublement
des jets qui conserve les $q$ occurrences d'une ligne de profondeur $q$ est
\begin{equation}\label{eq:pdf2-bsd-desdoblamiento-jets}
 \begin{aligned}
 \operatorname{Jet}_{\rm Boc,\mathbb Z}(E)
   &:=\bigoplus_{q\geq1}\bigoplus_{k=1}^{q}
       A_{q,\mathbb Z}\,e_{q,k},\\
 \mathcal U_{\rm jet}:
 \operatorname{Jet}_{\rm Boc,\mathbb Z}(E)
   &\xrightarrow{\ \sim\ }
     \bigoplus_{k\geq1}V_{k,\mathbb Z},\\
 \mathcal U_{\rm jet}((a_{q,k})_{q,k})
   &:=(v_k)_k,
   \qquad v_k:=\sum_{q\geq k}s_q(a_{q,k}).
 \end{aligned}
\end{equation}
L'inverse décompose chaque $v_k$ dans la somme directe précédente. En particulier,
\begin{equation}\label{eq:pdf2-bsd-dimension-jet}
 \operatorname{rank}_{\mathbb Z}\operatorname{Jet}_{\rm Boc,\mathbb Z}(E)
 =\sum_q q\,\operatorname{rank}_{\mathbb Z}A_{q,\mathbb Z}
 =\sum_k\operatorname{rank}_{\mathbb Z}V_{k,\mathbb Z}=d.
\end{equation}
Le recollement se restreint, dans ces mêmes bases, à
$\operatorname{gl}_{q,\mathbb Z}:V_{q,\mathbb Z}\to
\operatorname{Surv}_{{\rm det},\mathbb Z}^{\rm MW}$.

On définit, sans utiliser les dimensions,
\begin{equation}\label{eq:pdf2-bsd-psi-y-lambda-q}
 \begin{aligned}
  \psi_{q,\mathbb Z}&:=
      \operatorname{pg}_qP_E^{\rm PT}\kappa_{E,\mathbb Z}^{\rm HMT}:
      M_{E,\mathbb Z}^{\rm MW}\to V_{q,\mathbb Z},\\
  \lambda_{q,\mathbb Z}&:=
      \operatorname{pr}_{\rm MW}P_E^{\rm Man}
      \operatorname{gl}_{q,\mathbb Z}:
      V_{q,\mathbb Z}\to M_{E,\mathbb Z}^{\rm MW},\\
  \psi_q&:=\operatorname{pg}_qP_E^{\rm PT}\kappa_E^{\rm HMT}:
      M_E^{\rm MW}\to V_q,\\
  \lambda_q&:=\operatorname{pr}_{\rm MW}P_E^{\rm Man}
      \operatorname{gl}_q:V_q\to M_E^{\rm MW}.
 \end{aligned}
\end{equation}
Les domaines intermédiaires sont liés par les deux compatibilités basées
qui font partie du recollement Poitou--Tate et qui sont écrites ici sur son
action atteignable :
\begin{equation}\label{eq:pdf2-bsd-kummer-glue-compatibilidad}
 \kappa_E^{\rm HMT}\operatorname{pr}_{\rm MW}P_E^{\rm Man}
       \operatorname{gl}_q=\operatorname{gl}_q,
 \qquad
 \operatorname{pr}_{\rm MW}P_E^{\rm Man}\kappa_E^{\rm HMT}
       =I_{M_E^{\rm MW}}.
\end{equation}
Les deux mêmes identités, avec l'indice $\mathbb Z$, valent sur
$M_{E,\mathbb Z}^{\rm MW}$ et $V_{q,\mathbb Z}$, car toutes les matrices
précédentes sont intégrales et basées.

\begin{theorem}[Isomorphisme basé du drapeau Poitou--Tate]
\label{thm:pdf2-bsd-psi-pt}
Les applications
\begin{equation}\label{eq:pdf2-bsd-psi-pt}
 \begin{aligned}
 \Psi_{\rm PT}:M_E^{\rm MW}&\longrightarrow
       \operatorname{Flag}_{\rm Boc}(E):=\bigoplus_{q\ge1}V_q,
 &P&\longmapsto(\psi_q(P))_{q\ge1},\\
 \Psi_{\rm PT}^{-1}:\operatorname{Flag}_{\rm Boc}(E)&\longrightarrow M_E^{\rm MW},
 &(v_q)_q&\longmapsto\sum_{q\ge1}\lambda_q(v_q)
 \end{aligned}
\end{equation}
sont inverses. En conséquence,
\begin{equation}\label{eq:pdf2-bsd-psi-pt-integral}
 \Psi_{{\rm PT},\mathbb Z}:M_{E,\mathbb Z}^{\rm MW}
 \xrightarrow{\ \sim\ }\bigoplus_{q\ge1}V_{q,\mathbb Z},
 \qquad
 \Psi_{{\rm PT},\mathbb Z}^{-1}((v_q)_q)=
 \sum_q\lambda_{q,\mathbb Z}(v_q),
\end{equation}
et, seulement après extension des scalaires,
\begin{equation}\label{eq:pdf2-bsd-rango-orden}
 \operatorname{rank}E(\mathbb Q)
 =\sum_{q\ge1}\dim_K V_q
 =d=\operatorname{ord}_T\det S_E(T).
\end{equation}
\end{theorem}

\begin{proof}
Appliquer \eqref{eq:pdf2-bsd-unidad-counidad-paginas} avec
\eqref{eq:pdf2-bsd-kummer-glue-compatibilidad} donne, avec tous les domaines
intermédiaires visibles,
$\psi_r\lambda_q=\delta_{rq}I$ et
$\sum_q\lambda_q\psi_q=I_{M_E^{\rm MW}}$. Les deux formules de
\eqref{eq:pdf2-bsd-psi-pt} sont donc inverses avant de comparer les dimensions.
Les tables terminales sont intégrales et basées ; le même calcul, sans
dénominateurs, prouve \eqref{eq:pdf2-bsd-psi-pt-integral}. Ainsi, aucun indice de
sous-réseau n'apparaît lorsqu'on choisit ensuite une base de Mordell--Weil.
Ce n'est qu'après que l'on prend les dimensions et utilise
\eqref{eq:pdf2-bsd-orden-smith}. L'accouplement
$j_q\bar\beta_q$ transporte les bases saturées vers les lignes de hauteur dont la
comparaison avec Néron--Tate s'écrit dans
\eqref{eq:pdf2-bsd-altura-nt-generadores}.
\end{proof}

\section{Localisation, rang de Smith et exactitude Sha--Tamagawa}

Dans cette section, $v$ parcourt les places finies ; aux places de bonne réduction, on
a $\Phi_v(\mathbb F_v)=0$ et $c_v=1$, de sorte que les sommes et produits
locaux ont un support fini. La place infinie est réservée à l'évaluation
de la différentielle de Néron dans la quatrième flèche. Pour chaque place finie $v$, on
a le triangle
\[
 \mathcal C_{E,v}^{\rm fin}\to\mathcal C_{E,v}\to
 \mathcal C_{E,v}^{\rm sing}\to\mathcal C_{E,v}^{\rm fin}[1]
\]
et l'application de chaînes
\begin{equation}\label{eq:pdf2-bsd-ell-e}
 \ell_E(x,(y_v)_v)=(\operatorname{res}_vx-i_vy_v)_v.
\end{equation}
Le cône réduit est
\begin{equation}\label{eq:pdf2-bsd-cono-reducido}
 \mathcal C_E^{\rm loc,red}:=
 q_{\rm root}\operatorname{Cone}(\ell_E)[-1].
\end{equation}
En degrés de cocycles, un élément s'écrit $(x,(y_v),(h_v))$ et
\begin{equation}\label{eq:pdf2-bsd-diferencial-cono}
 d(x,(y_v),(h_v))=
 (dx,(dy_v),(\operatorname{res}_vx-i_vy_v-dh_v)_v).
\end{equation}

Après avoir séparé par \eqref{eq:pdf2-bsd-psi-pt} le réseau libre de Mordell--Weil
et son dual, les blocs FS sont contractés par
\eqref{eq:pdf2-bsd-hfs-global}. Dans le réseau restant, l'inverse
\eqref{eq:pdf2-bsd-j-e-inverso-racional} et ses deux vérifications
\eqref{eq:pdf2-bsd-j-e-dos-lados} donnent, sans déduire l'acyclicité de la seule
dualité,
\begin{equation}\label{eq:pdf2-bsd-rango-completo}
 D_E\otimes\mathbb Q:M_E\otimes\mathbb Q
 \xrightarrow{\sim}N_E\otimes\mathbb Q.
\end{equation}
La forme normale intégrale donne
\begin{equation}\label{eq:pdf2-bsd-smith-finito}
 UD_EV=\operatorname{diag}(s_1,\ldots,s_t),
 \qquad
 F_E:=\operatorname{coker}D_E,
 \qquad |F_E|=\prod_{i=1}^t|s_i|<\infty.
\end{equation}

Soit
\begin{equation}\label{eq:pdf2-bsd-pi-componentes}
 \pi_\Phi:N_E\longrightarrow\bigoplus_v\Phi_v(\mathbb F_v)
\end{equation}
la réduction de chaque coordonnée locale modulo
$\mathcal C_{E,v}^{\rm fin}+d\mathcal C_{E,v}$ ; on a
$\pi_\Phi D_E=0$. Pour une classe
$[\xi]\in\operatorname{Sha}(E)$, choisissons ses primitives finies
$y_v$ et des homotopies $h_v$ avec
$\operatorname{res}_v\xi-i_vy_v=dh_v$. La formule
\begin{equation}\label{eq:pdf2-bsd-iota-sha-cociclos}
 \begin{aligned}
 n_\xi&:=\operatorname{pr}_{N_E}\operatorname{pr}_{\mathcal L_E}
       \Theta_E(\xi,(y_v),(h_v)),\\
 \iota_{\rm Sha}[\xi]&:=[n_\xi]\in F_E
 \end{aligned}
\end{equation}
est indépendante des choix par
\eqref{eq:pdf2-bsd-diferencial-cono}.

Pour rendre matérielle la surjectivité finale, soit
$\widetilde\phi_v$ le représentant terminal basé de
$\phi_v\in\Phi_v(\mathbb F_v)$. La somme des counités est corrigée dans
l'unité commune et relevée par l'opérateur structurel :
\begin{equation}\label{eq:pdf2-bsd-lift-componentes}
 \operatorname{Lift}_\Phi((\phi_v)_v):=
 \operatorname{pr}_{N_E}\operatorname{Ext}_{\Gamma,{\rm det}}
 \left((\widetilde\phi_v)_v,
 -\widehat\Pi_0\!\left(\sum_v
       \varepsilon_{\rm AW}(\widetilde\phi_v)\right)\right).
\end{equation}
La relation de corestriction de
$\operatorname{Rel}_{\rm det}$ fait de cette expression un cocycle et donne
\begin{equation}\label{eq:pdf2-bsd-lift-dos-identidades}
 \pi_\Phi\operatorname{Lift}_\Phi=I,
 \qquad
 \operatorname{im}D_E=
 \ker\pi_\Phi\cap\ker(\text{classe globale}).
\end{equation}

\begin{theorem}[Suite exacte finie Sha--Tamagawa]
\label{thm:pdf2-bsd-sha-tamagawa}
Les applications de cocycles précédentes induisent
\begin{equation}\label{eq:pdf2-bsd-sha-tamagawa}
 0\longrightarrow\operatorname{Sha}(E)
 \xrightarrow{\ \iota_{\rm Sha}\ }F_E
 \xrightarrow{\ \partial_{\rm loc}\ }
 \bigoplus_v\Phi_v(\mathbb F_v)
 \longrightarrow0,
 \qquad
 \partial_{\rm loc}[n]:=\pi_\Phi(n).
\end{equation}
En particulier, $\operatorname{Sha}(E)$ est fini et
\begin{equation}\label{eq:pdf2-bsd-cardinal-sts}
 |F_E|=|\operatorname{Sha}(E)|\prod_vc_v,
 \qquad c_v:=|\Phi_v(\mathbb F_v)|.
\end{equation}
\end{theorem}

\begin{proof}
L'égalité $\pi_\Phi D_E=0$ rend bien définie
$\partial_{\rm loc}$. Si \eqref{eq:pdf2-bsd-iota-sha-cociclos} est un
bord, sa composante globale est un bord ; ainsi, $\iota_{\rm Sha}$ est
injective. Une classe de $F_E$ est dans $\ker\partial_{\rm loc}$ si et seulement
si toutes ses coordonnées singulières admettent des représentants finis ; par
\eqref{eq:pdf2-bsd-diferencial-cono}, cela équivaut à appartenir à l'image de
$\iota_{\rm Sha}$. Enfin,
\eqref{eq:pdf2-bsd-lift-dos-identidades} donne un antécédent de chaque tuple de
composantes et prouve la surjectivité. Les termes adjacents libres de
la suite longue sont précisément les deux termes séparés par
$\Psi_{\rm PT}$ et $X^{\rm orth}$ ; les termes adjacents FS sont
contractiles par $H_{\rm FS}$. Leur annulation n'a pas été supposée. La finitude
de $F_E$ provient de \eqref{eq:pdf2-bsd-rango-completo} ; l'injection donne la
 finitude de $\operatorname{Sha}(E)$ et prendre les ordres donne
 \eqref{eq:pdf2-bsd-cardinal-sts}.
\end{proof}

\begin{corollary}[Satisfaction primaire complète et recomposition intégrale]
\label{cor:pdf2-bsd-sha-tamagawa-todos-ell}
Pour chaque nombre premier $\ell$, le changement de base des applications de cocycles produit
la suite exacte
\begin{equation}\label{eq:pdf2-bsd-sha-tamagawa-ell}
 0\longrightarrow\operatorname{Sha}(E)[\ell^\infty]
 \xrightarrow{\ \iota_{{\rm Sha},\ell}\ }F_{E,\ell}
 \xrightarrow{\ \partial_{{\rm loc},\ell}\ }
 \bigoplus_v\Phi_v(\mathbb F_v)[\ell^\infty]
 \longrightarrow0,
\end{equation}
où
\begin{equation}\label{eq:pdf2-bsd-mapas-ell-generadores}
 \iota_{{\rm Sha},\ell}([\xi]\otimes1)
   =[n_\xi\otimes1],
 \qquad
 \partial_{{\rm loc},\ell}([n\otimes1])
   =(\pi_{\Phi,v}(n)\otimes1)_v.
\end{equation}
De plus, la somme de toutes les suites
\eqref{eq:pdf2-bsd-sha-tamagawa-ell} est exactement
\eqref{eq:pdf2-bsd-sha-tamagawa} ; ce n'est pas seulement sa partie $3$-primaire.
\end{corollary}

\begin{proof}
La matrice $D_E$, le cocycle $n_\xi$, la réduction $\pi_\Phi$ et
$\operatorname{Lift}_\Phi$ sont intégraux par
\eqref{eq:pdf2-bsd-complejos-base-change}--
\eqref{eq:pdf2-bsd-theta-base-change}. Comme $\mathbb Z_\ell$ est plat sur
$\mathbb Z$, tensoriser la suite exacte finie
\eqref{eq:pdf2-bsd-sha-tamagawa} conserve l'exactitude. Pour un groupe fini
$A$, on a canoniquement
$A\otimes\mathbb Z_\ell=A[\ell^\infty]$ ; cela identifie les trois termes
et donne les formules sur les générateurs
\eqref{eq:pdf2-bsd-mapas-ell-generadores}. La forme de Smith
\eqref{eq:pdf2-bsd-smith-finito} écrit
$s_i=\prod_\ell\ell^{v_\ell(s_i)}$, de sorte que
\[
 F_E\cong\bigoplus_\ell F_{E,\ell},\qquad
 \operatorname{Sha}(E)\cong
   \bigoplus_\ell\operatorname{Sha}(E)[\ell^\infty],\qquad
 \Phi_v(\mathbb F_v)\cong
   \bigoplus_\ell\Phi_v(\mathbb F_v)[\ell^\infty].
\]
Seuls un nombre fini de nombres premiers et un nombre fini de places interviennent car tous les groupes
sont des quotients du groupe fini $F_E$. Sommer les suites primaires restitue
la suite intégrale. En prenant les ordres d'abord premier par premier puis en multipliant,
\begin{equation}\label{eq:pdf2-bsd-cardinal-sts-todos-ell}
 \prod_\ell|F_{E,\ell}|
 =\prod_\ell|\operatorname{Sha}(E)[\ell^\infty]|
   \prod_v\prod_\ell|\Phi_v(\mathbb F_v)[\ell^\infty]|
 =|\operatorname{Sha}(E)|\prod_vc_v.
\end{equation}
Ainsi, tous les coefficients et tous les
facteurs de Tamagawa sont matériellement traités.
\end{proof}

\section{Comparateur analytique--arithmétique basé}

La coordonnée analytique n'est pas postulée à partir de l'ordre recherché. Soit
$\lambda_3=\log4$, correspondant au générateur cyclotomique basé
$\gamma_3=4$, et définissons dans un disque $U_E$ centré en $1$
\begin{equation}\label{eq:pdf2-bsd-t-e-construido}
 T_E(s):=\frac{\exp(\lambda_3(s-1))-1}{\lambda_3}
 =(s-1)+\frac{\lambda_3}{2}(s-1)^2+O((s-1)^3).
\end{equation}
Par calcul direct,
\begin{equation}\label{eq:pdf2-bsd-t-e-normalizacion-probada}
 T_E(1)=0,
 \qquad T_E'(s)=\exp(\lambda_3(s-1)),
 \qquad T_E'(1)=1.
\end{equation}
Le théorème d'inversion locale fournit, après réduction de $U_E$ si
nécessaire, une carte biholomorphe $s\leftrightarrow T_E(s)$.
Pour comparer toutes les fibres, prenons
$\gamma_\ell=1+\ell$ si $\ell$ est impair et $\gamma_2=5$, et posons
\begin{equation}\label{eq:pdf2-bsd-t-ell-construido}
 T_{E,\ell}(s):=
 \frac{\exp(\log(\gamma_\ell)(s-1))-1}{\log(\gamma_\ell)}.
\end{equation}
Les cartes sont reliées par un unique automorphisme formel
$T_{E,\ell}=\vartheta_{\ell,3}(T_E)$ avec
$\vartheta_{\ell,3}(0)=0$ et $\vartheta_{\ell,3}'(0)=1$. Ainsi, ni
l'ordre en $T$ ni le coefficient principal ne changent en passant de $3$ à $\ell$.

Le prolongement analytique qui atteint cette carte n'est pas non plus déclaré. Par le
théorème de modularité, soit
\[
 f_E(z)=\sum_{m\geq1}a_m(E)e^{2\pi imz}
\]
la forme nouvelle normalisée de poids $2$ et de niveau $N_E$ associée à $E$, et soit
$\varepsilon_E\in\{\pm1\}$ le signe fixé par
\begin{equation}\label{eq:pdf2-bsd-transformacion-modular}
 g_E(y):=f_E\!\left(\frac{iy}{\sqrt{N_E}}\right),
 \qquad
 g_E(1/y)=\varepsilon_Ey^2g_E(y).
\end{equation}
Pour $\Re s>3/2$, l'intégration terme à terme donne
\begin{equation}\label{eq:pdf2-bsd-mellin-semiplano}
 \Lambda_E(s):=N_E^{s/2}(2\pi)^{-s}\Gamma(s)L(E,s)
 =\int_0^\infty g_E(y)y^s\,\frac{dy}{y}.
\end{equation}
On partage l'intégrale en $(0,1)$ et $(1,\infty)$ et, dans la première partie, on
effectue le changement $y=1/t$. La deuxième égalité de
\eqref{eq:pdf2-bsd-transformacion-modular} produit la formule matérielle
\begin{equation}\label{eq:pdf2-bsd-mellin-continuada}
 \Lambda_E^{\rm cont}(s)=
 \int_1^\infty g_E(y)
       \bigl(y^s+\varepsilon_Ey^{2-s}\bigr)\,\frac{dy}{y}.
\end{equation}
La fonction $g_E(y)$ décroît exponentiellement lorsque $y\to\infty$ ;
l'intégrale de \eqref{eq:pdf2-bsd-mellin-continuada} et toutes ses dérivées en
$s$ convergent donc uniformément sur les compacts. Elle définit ainsi une fonction entière et
coïncide avec \eqref{eq:pdf2-bsd-mellin-semiplano} là où la série de Dirichlet
converge. Puisque
$N_E^{s/2}(2\pi)^{-s}\Gamma(s)$ est holomorphe et non nulle près de $s=1$, la
division par ce facteur construit, sans utiliser son ordre d'annulation, le germe
holomorphe de $L(E,s)$ dans un disque $U_E$ centré en $1$.

Pour éviter de mélanger une fibre $\ell$-adique avec une matrice complexe, on prend
d'abord le polynôme entier, indépendant de $\ell$,
\begin{equation}\label{eq:pdf2-bsd-polinomio-local-entero}
 P_{E,p}(X):=\det(I-X\operatorname{Frob}_p\mid V_\ell E^{I_p})
 \in\mathbb Z[X].
\end{equation}
Pour $p\nmid N_E$, soient $\alpha_p,\beta_p$ les racines de $P_{E,p}$. On prend
$W_p=\mathbb C^2$ avec une base orthonormale et la réalisation normale
\begin{equation}\label{eq:pdf2-bsd-realizacion-normal-frobenius}
 F_p^{\mathbb C}:=\operatorname{diag}(\alpha_p,\beta_p),
 \qquad |\alpha_p|=|\beta_p|=p^{1/2},
 \qquad \|F_p^{\mathbb C}\|_1=2p^{1/2}.
\end{equation}
Pour les mauvais nombres premiers, on choisit de même une réalisation complexe du
polynôme local ; ils sont en nombre fini et sont conservés comme blocs séparés. On n'utilise pas
la matrice compagnon, dont la norme ne serait pas contrôlée uniformément par le
module de ses valeurs propres.

Posons $\mathscr H_E:=\widehat\bigoplus_{p\nmid N_E}W_p$ et, pour $X>0$,
\begin{equation}\label{eq:pdf2-bsd-truncados-fredholm}
 K_{E,X}(s):=\bigoplus_{\substack{p\leq X\\p\nmid N_E}}
          p^{-s}F_p^{\mathbb C},
 \qquad
 L_X(E,s):=\det(I-K_{E,X}(s))^{-1}
       \prod_{p\mid N_E}P_{E,p}(p^{-s})^{-1}.
\end{equation}
Si $K_0$ est un compact de $\{\Re s>3/2\}$ et
$\sigma_0:=\inf_{s\in K_0}\Re s$, alors
\begin{equation}\label{eq:pdf2-bsd-cota-traza}
 \sup_{s\in K_0}\|K_{E,X'}(s)-K_{E,X}(s)\|_1
 \leq C_E\sum_{X<p\leq X'}p^{1/2-\sigma_0}
 \xrightarrow[X,X'\to\infty]{}0.
\end{equation}
Ainsi, $K_{E,X}\to K_E$ en norme trace, localement uniformément, et la
continuité du déterminant de Fredholm prouve, et ne définit pas,
\begin{equation}\label{eq:pdf2-bsd-fredholm-euler}
 \det_F(I-K_E(s))^{-1}
 \prod_{p\mid N_E}P_{E,p}(p^{-s})^{-1}
 =\lim_{X\to\infty}L_X(E,s)=L(E,s).
\end{equation}

Pour passer du produit infini à un complexe fini, on conserve d'abord
l'opérateur, non seulement son déterminant. Soient $\mathscr H_E^+$ et
$\mathscr H_E^-$ deux copies basées de
$\mathscr H_E\oplus\bigoplus_{p\mid N_E}W_p$. Dans le demi-plan
$\Re s>3/2$, on définit
\begin{equation}\label{eq:pdf2-bsd-operadores-euler-inversos}
 \begin{aligned}
  \mathscr A_{E,X}(s)&:=
   (I-K_{E,X}(s))^{-1}\oplus
   \bigoplus_{p\mid N_E}(I-p^{-s}F_p^{\mathbb C})^{-1},\\
  \mathscr A_E(s)&:=\lim_{X\to\infty}\mathscr A_{E,X}(s):
       \mathscr H_E^+\longrightarrow\mathscr H_E^-.
 \end{aligned}
\end{equation}
L'identité de la résolvante et \eqref{eq:pdf2-bsd-cota-traza} donnent, pour chaque
compact $K_0$ du demi-plan,
\begin{equation}\label{eq:pdf2-bsd-resolvente-convergencia-traza}
 \sup_{s\in K_0}
 \|\mathscr A_{E,X'}(s)-\mathscr A_{E,X}(s)\|_1\longrightarrow0.
\end{equation}
Soit $\mathcal U_{\rm bas}:\mathscr H_E^+\to\mathscr H_E^-$ l'identification unitaire qui
envoie chaque vecteur de la base positive sur le vecteur de même nom de la base
négative. Les déterminants suivants désignent ceux de
$\mathcal U_{\rm bas}^{-1}\mathscr A_{E,X}=I+\mathcal L^1$ et
$\mathcal U_{\rm bas}^{-1}\mathscr A_E=I+\mathcal L^1$. Alors
\begin{equation}\label{eq:pdf2-bsd-det-a-x-l-x}
 \det_F(\mathcal U_{\rm bas}^{-1}\mathscr A_{E,X}(s))=L_X(E,s),
 \qquad
 \det_F(\mathcal U_{\rm bas}^{-1}\mathscr A_E(s))=L(E,s).
\end{equation}
La deuxième égalité est précisément la limite de
\eqref{eq:pdf2-bsd-fredholm-euler} ; le signe est correct car on prend
l'inverse de $I-K_E$.

La corestriction finie se construit à partir des terminaux, sans publier
une infinité de nombres premiers dans une matrice finie. Soient $\mathcal I_E^+$ et
$\mathcal I_E^-$ les listes terminales des bases intégrales de $M_E$ et
$N_E$, respectivement ; elles sont finies par
\eqref{eq:pdf2-bsd-sumando-fs-y-lattice}. $\sigma_*>3/2$ étant fixé,
l'entrelaceur local--global les insère dans l'espace d'Euler par
\begin{equation}\label{eq:pdf2-bsd-embedding-terminal-euler}
 \iota_{\rm Eul}^\pm(e_\eta):=
 \left(p^{-\sigma_*}\operatorname{res}_p^\pm
       \operatorname{Surv}_E(e_\eta)\right)_p
 \oplus\left(\operatorname{res}_p^\pm
       \operatorname{Surv}_E(e_\eta)\right)_{p\mid N_E}.
\end{equation}
La bonne somme est de carré sommable par la même borne que
\eqref{eq:pdf2-bsd-cota-traza}. On pose
$\tau_\eta^\pm:=\iota_{\rm Eul}^\pm(e_\eta)$ ; l'indépendance des lignes
terminales donne les counités duales $\{\epsilon_{\eta}^{\pm}\}$, de sorte que
$\epsilon_\eta^\pm(\tau_{\eta'}^\pm)=\delta_{\eta\eta'}$. Les sommes
\begin{equation}\label{eq:pdf2-bsd-proyecciones-terminales-finitas}
 P_E^\pm:=\sum_{\eta\in\mathcal I_E^\pm}
      \tau_\eta^\pm\epsilon_\eta^\pm,
 \qquad Q_E^\pm:=I-P_E^\pm
\end{equation}
sont des projections basées de rang fini ; leurs indices sont fixés par les
étiquettes terminales et non par le rang de Mordell--Weil ni par un ordre
d'annulation.

La coordonnée $\deg$ du ledger décompose
\begin{equation}\label{eq:pdf2-bsd-graduacion-fredholm}
 \mathscr H_E^\pm=\widehat\bigoplus_i\mathscr H_E^{\pm,i},
 \qquad
 \mathcal I_E^{\pm,i}:=\{\eta\in\mathcal I_E^\pm:\deg\eta=i\},
 \qquad
 P_E^{\pm,i}:=
 \sum_{\eta\in\mathcal I_E^{\pm,i}}\tau_\eta^\pm\epsilon_\eta^\pm.
\end{equation}
Les relations locales conservent $\deg$ sauf la différentielle $i\mapsto i+1$ ;
par conséquent,
$\mathscr A_E=\bigoplus_i\mathscr A_E^i$ avec
$\mathscr A_E^i:\mathscr H_E^{+,i}\to\mathscr H_E^{-,i+1}$. Les quatre
blocs de \eqref{eq:pdf2-bsd-bloques-fredholm}, son Schur et ses troncatures
sont pris degré par degré, et, par définition,
\begin{equation}\label{eq:pdf2-bsd-schur-graduado}
 \mathscr D_E=\bigoplus_i\mathscr D_E^i,
 \qquad
 \mathscr D_E^i:=A_{00}^i-A_{01}^i(A_{11}^i)^{-1}A_{10}^i.
\end{equation}

Le prolongement de Fredholm est écrit avant son utilisation. Dans la carte d'Euler
$U_0$, on pose $\mathscr A_{E,0}:=\mathscr A_E$. Les cartes conservent un
complément basé commun $\mathscr Q_E^\pm$ et ne changent que la famille finie
des terminaux. Si
$\tau_{\eta,\alpha}^\pm:=
\operatorname{Ext}_{\Gamma,{\rm det},\alpha}(\tau_\eta^\pm)$ et
$\epsilon_{\eta,\alpha}^\pm$ est sa base duale, on définit
\begin{equation}\label{eq:pdf2-bsd-cambios-fredholm-cartas}
 C_{\alpha\beta}^\pm:=
 I_{\mathscr Q_E^\pm}\oplus
 \sum_{\eta\in\mathcal I_E^\pm}
 \tau_{\eta,\beta}^\pm\epsilon_{\eta,\alpha}^\pm,
 \qquad C_{\alpha\beta}^\pm-I\in\mathcal L^1,
 \qquad
 C_{\beta\gamma}^\pm C_{\alpha\beta}^\pm=C_{\alpha\gamma}^\pm.
\end{equation}
La différence par rapport à l'identité a un rang au plus égal à
$|\mathcal I_E^\pm|$ et est donc de classe trace ; l'holomorphie découle
du fait que chacune des coordonnées terminales, en nombre fini, de
$\operatorname{Ext}_{\Gamma,{\rm det},\alpha}$ est holomorphe. La dernière
égalité se vérifie sur $\mathscr Q_E^\pm$ et sur chaque
$\tau_{\eta,\alpha}^\pm$. On définit
récursivement sur les recouvrements
\begin{equation}\label{eq:pdf2-bsd-continuacion-fredholm-cartas}
 \mathscr A_{E,\beta}:=
 C_{\alpha\beta}^-\mathscr A_{E,\alpha}
 (C_{\alpha\beta}^+)^{-1}.
\end{equation}
Les opérateurs $C_{\alpha\beta}^\pm$ sont holomorphes car chaque coefficient
est une coordonnée de $\operatorname{Ext}_{\Gamma,{\rm det}}$ ; le cocycle de
\eqref{eq:pdf2-bsd-cambios-fredholm-cartas} prouve que
\eqref{eq:pdf2-bsd-continuacion-fredholm-cartas} ne dépend pas de la chaîne de
cartes. C'est la famille de classe déterminant désormais notée
$\mathscr A_E(s)$. Avec
$C_{0\alpha}^\pm:=C_{\alpha-1,\alpha}^\pm\cdots C_{0,1}^\pm$, on transporte
en outre
$P_{E,\alpha}^{\pm,i}:=
C_{0\alpha}^\pm P_E^{\pm,i}(C_{0\alpha}^\pm)^{-1}$ et
$\mathscr H_{E,\alpha}^{\pm,i}:=
C_{0\alpha}^\pm\mathscr H_E^{\pm,i}$ ; ainsi, tous les Schur
$\mathscr D_{E,\alpha}^i$ sont ceux de
\eqref{eq:pdf2-bsd-schur-graduado}, non de nouveaux symboles. Après réduction de chaque
carte, le bloc
\begin{equation}\label{eq:pdf2-bsd-bloques-fredholm}
 A_{ab}(s):=P_a^-\mathscr A_E(s)P_b^+,
 \qquad P_0^\pm:=P_E^\pm,\quad P_1^\pm:=Q_E^\pm,
 \qquad A_{11}(s):Q_E^+\mathscr H_E^+\xrightarrow{\sim}
                       Q_E^-\mathscr H_E^-
\end{equation}
est inversible. C'est la réduction finie standard de l'opérateur de Fredholm ;
le choix concret de $P_E^\pm$ est donné par
\eqref{eq:pdf2-bsd-proyecciones-terminales-finitas}. La base fixe aussi
l'identification unitaire
$U_{11}:Q_E^+\mathscr H_E^+\to Q_E^-\mathscr H_E^-$ ; par construction,
$U_{11}^{-1}A_{11}=I+\mathcal L^1$, et c'est seulement à cet opérateur que s'applique
$\det_F$. Les copies $+$ et $-$ sont déjà la totalisation paire--impaire ; on
n'applique pas un deuxième signe alterné au déterminant de l'opérateur inverse.

Pour chaque troncature qui contient le support de ces terminaux, on définit
matériellement
\begin{equation}\label{eq:pdf2-bsd-cor-lift-truncados}
 \begin{aligned}
  \operatorname{Cor}_{E,X}&:=P_E^-,\\
  \operatorname{Lift}_{E,X}(v)&:=
       v-A_{11,X}^{-1}A_{10,X}v,\\
  \mathscr D_{E,X}(s)&:=
       \operatorname{Cor}_{E,X}\mathscr A_{E,X}(s)
       \operatorname{Lift}_{E,X}
       =A_{00,X}-A_{01,X}A_{11,X}^{-1}A_{10,X}.
 \end{aligned}
\end{equation}
Par multiplication de blocs,
\begin{equation}\label{eq:pdf2-bsd-eliminacion-schur}
 \begin{pmatrix}I&-A_{01,X}A_{11,X}^{-1}\\0&I\end{pmatrix}
 \mathscr A_{E,X}
 \begin{pmatrix}I&0\\-A_{11,X}^{-1}A_{10,X}&I\end{pmatrix}
 =\begin{pmatrix}\mathscr D_{E,X}&0\\0&A_{11,X}\end{pmatrix}.
\end{equation}
Ainsi, la corestriction n'est pas nominale : son relèvement, son inverse sur le complément
et sa différentielle finie sont écrits. La convergence de la résolvante
implique
\begin{equation}\label{eq:pdf2-bsd-corestriccion-limite}
 \mathscr D_E(s):=\lim_{X\to\infty}\mathscr D_{E,X}(s)
 =A_{00}-A_{01}A_{11}^{-1}A_{10},
 \qquad
 \|\mathscr D_{E,X}-\mathscr D_E\|_1\to0.
\end{equation}
Le déterminant du complément n'est pas écarté : par
\eqref{eq:pdf2-bsd-eliminacion-schur},
\begin{equation}\label{eq:pdf2-bsd-determinante-schur-complemento}
 \det_F(\mathcal U_{\rm bas}^{-1}\mathscr A_E)
 =\det_F(U_{11}^{-1}A_{11})\det(\mathscr D_E).
\end{equation}

Soit $\mathcal O(U_E)$ l'anneau des fonctions holomorphes du disque central.
Les espaces rationnels basés $V_0^i,W_0^{i+1}$ de
\eqref{eq:pdf2-bsd-s-universal} fournissent
\begin{equation}\label{eq:pdf2-bsd-dominios-analiticos}
 \mathscr V_E^i:=V_0^i\otimes_{\mathbb Q}\mathcal O(U_E),
 \qquad
 \mathscr W_E^{i+1}:=W_0^{i+1}\otimes_{\mathbb Q}\mathcal O(U_E).
\end{equation}
Les lignes de $\operatorname{Rel}_{\rm det}$ construisent, sur chaque terminal
$\eta$ et à chaque degré, les vecteurs du modèle commun
$v_{\eta,i}:=\upsilon_i(b_\eta)$,
$w_{\eta,i+1}:=\upsilon_{i+1}(b_\eta^{\rm term})$ et les isomorphismes basés
\begin{equation}\label{eq:pdf2-bsd-jvw-generadores}
 \begin{aligned}
  J_{E,V,i}(\tau_{\eta,i}^+)&:=v_{\eta,i},\\
  J_{E,W,i+1}(\tau_{\eta,i+1}^-)&:=w_{\eta,i+1}.
 \end{aligned}
\end{equation}
La relation différentielle de cette même ligne, appliquée à chaque générateur, est
\begin{equation}\label{eq:pdf2-bsd-entrelazamiento-analitico-smith}
 \boxed{
 J_{E,W,i+1}(s)\,\mathscr D_E^i(s)
 =S_E^i(T_E(s))\,J_{E,V,i}(s).}
\end{equation}
C'est l'entrelaceur effectif entre le Schur de Fredholm corestreint et la
matrice de Smith ; ce n'est pas une égalité de déterminants postulée.

Pour typer le prolongement, soit $\mathfrak U_E$ un domaine connexe qui contient
un ouvert du demi-plan $\Re s>3/2$ et le disque $U_E$, et choisissons une chaîne
finie de cartes $\{U_\alpha\}_{\alpha=0}^A$ de
$\operatorname{Ext}_{\Gamma,{\rm det}}$, avec $U_A=U_E$. Dans $U_\alpha$, on
pose
$\mathscr P_{E,\alpha}^i:=
P_{E,\alpha}^{+,i}\mathscr H_{E,\alpha}^{+,i}$ et
$\mathscr D_{E,\alpha}^i$ égal au Schur de
\eqref{eq:pdf2-bsd-corestriccion-limite}. Les applications vers le modèle commun sont
définies, sans choix supplémentaire, par
\begin{equation}\label{eq:pdf2-bsd-j-cartas-modelo-comun}
 J_{E,\alpha,i}(\tau_{\eta,\alpha,i}):=v_{\eta,i}
 \quad\text{dans le domaine},
 \qquad
 J_{E,\alpha,i+1}(\tau_{\eta,\alpha,i+1}):=w_{\eta,i+1}
 \quad\text{dans le codomaine}.
\end{equation}
Ils sont holomorphes car les bases de
\eqref{eq:pdf2-bsd-cambios-fredholm-cartas} le sont. Sur un recouvrement, on
définit, toujours au même degré,
\begin{equation}\label{eq:pdf2-bsd-transiciones-atlas}
 G_{\alpha\beta,i}:=J_{E,\beta,i}^{-1}J_{E,\alpha,i}:
 \mathscr P_{E,\alpha}^i\big|_{U_\alpha\cap U_\beta}
 \longrightarrow
 \mathscr P_{E,\beta}^i\big|_{U_\alpha\cap U_\beta}.
\end{equation}
La composition littérale de ces matrices prouve
\begin{equation}\label{eq:pdf2-bsd-cociclo-atlas}
 G_{\beta\gamma,i}G_{\alpha\beta,i}=G_{\alpha\gamma,i},
 \qquad G_{\alpha\alpha,i}=I,
 \qquad G_{\beta\alpha,i}=G_{\alpha\beta,i}^{-1},
\end{equation}
et l'entrelacement dans les deux cartes donne la compatibilité de chaîne
\begin{equation}\label{eq:pdf2-bsd-cociclo-cadena}
 G_{\alpha\beta,i+1}\mathscr D_{E,\alpha}^i
 =\mathscr D_{E,\beta}^iG_{\alpha\beta,i}.
\end{equation}
Par conséquent,
$g_{\alpha\beta}:=\prod_i\det(G_{\alpha\beta,i})^{(-1)^i}$ satisfait le
cocycle triple. Si
\begin{equation}\label{eq:pdf2-bsd-tau-atlas}
 \tau_E^{\rm atlas}(s):=
 \prod_{\alpha=0}^{A-1}g_{\alpha,\alpha+1}(s),
\end{equation}
alors $\tau_E^{\rm atlas}$ est exactement le transport déterminantal
de la carte d'Euler vers la carte centrale.

On note $J_{E,V}$ et $J_{E,W}$ les totalisations paires--impaires, avec
l'ordre de Koszul, des applications de
\eqref{eq:pdf2-bsd-entrelazamiento-analitico-smith} ; leurs déterminants sont
les produits alternés des déterminants par degré. Dans la carte centrale,
on définit, et pas avant, l'unité conjointe
\begin{equation}\label{eq:pdf2-bsd-u-e-construida}
 u_E(s):=\tau_E^{\rm atlas}(s)
 \det_F(U_{11}^{-1}A_{11}(s))
 \frac{\det J_{E,V}(s)}
      {\det J_{E,W}(s)}.
\end{equation}
Tous ses facteurs sont holomorphes et inversibles en réduisant $U_E$. La
normalisation n'est pas imposée facteur par facteur. Soit $\mathbf e_{\rm cent}$ le
générateur de la ligne déterminante du modèle commun. En effectuant les deux
opérations triangulaires de \eqref{eq:pdf2-bsd-eliminacion-schur}, en transportant
par \eqref{eq:pdf2-bsd-cambios-fredholm-cartas} et en appliquant
\eqref{eq:pdf2-bsd-j-cartas-modelo-comun}, l'action conjointe sur $s=1$ est,
générateur par générateur,
\begin{equation}\label{eq:pdf2-bsd-u-e-normalizacion-probada}
 \begin{aligned}
 &\tau_E^{\rm atlas}(1)
 \det_F(U_{11}^{-1}A_{11}(1))
 \frac{\det J_{E,V}(1)}
      {\det J_{E,W}(1)}\,\mathbf e_{\rm cent}\\
 &\qquad=
 \bigotimes_{i,\eta}
 \bigl(\epsilon_{\eta,i}^-(\tau_{\eta,i}^-)
       \epsilon_{\eta,i}^+(\tau_{\eta,i}^+)^{-1}\bigr)^{(-1)^i}
 \mathbf e_{\rm cent}
 =\mathbf e_{\rm cent},
 \qquad \boxed{u_E(1)=1}.
 \end{aligned}
\end{equation}
La première égalité est l'identité de conservation de la base commune du
comparateur de $\mathfrak G_{\rm det}(E)$, maintenant développée par ses lignes ; la
deuxième utilise
$\epsilon_{\eta,i}^\pm(\tau_{\eta,i}^\pm)=1$. Aucun
des quatre facteurs n'a été normalisé séparément et la valeur recherchée n'a pas été utilisée.

\begin{proposition}[Provenance FERS de la factorisation analytique]
\label{prop:pdf2-bsd-fers-factorizacion-owner}
Dans le disque $U_E$, on a l'identité de sections déterminantales
\begin{equation}\label{eq:pdf2-bsd-factorizacion-analitica}
 \boxed{L(E,s)=u_E(s)\det S_E(T_E(s)).}
\end{equation}
L'égalité est la publication du lecteur FERS de
$\mathfrak G_{\rm det}(E)$ et non une égalité postulée après
le calcul de l'ordre ou du terme principal.
\end{proposition}

\begin{proof}
Dans le demi-plan absolu,
\eqref{eq:pdf2-bsd-det-a-x-l-x} identifie la limite de Fredholm à $L(E,s)$.
L'élimination littérale \eqref{eq:pdf2-bsd-eliminacion-schur} et la continuité
en norme trace donnent
$L(E,s)=\det_F(U_{11}^{-1}A_{11}(s))\det\mathscr D_E(s)$.
Prendre les produits extérieurs ordonnés dans
\eqref{eq:pdf2-bsd-entrelazamiento-analitico-smith} donne
\[
 \det\mathscr D_E(s)
 =\frac{\det J_{E,V}(s)}
        {\det J_{E,W}(s)}
   \det S_E(T_E(s)).
\]
Les équations \eqref{eq:pdf2-bsd-cociclo-atlas} et
\eqref{eq:pdf2-bsd-cociclo-cadena} transportent cette section à travers les cartes ;
le produit exact de leurs transitions est $\tau_E^{\rm atlas}$. Par conséquent, dans
$U_A=U_E$, le coefficient qui reste est précisément
\eqref{eq:pdf2-bsd-u-e-construida}, ce qui prouve
\eqref{eq:pdf2-bsd-factorizacion-analitica}. L'unicité du prolongement de
Mellin de \eqref{eq:pdf2-bsd-mellin-continuada} exclut une autre section. Enfin,
\eqref{eq:pdf2-bsd-u-e-normalizacion-probada} prouve conjointement
$u_E(1)=1$. Aucune étape n'utilise $d$, le rang, $\operatorname{Sha}(E)$ ni le
terme principal.
\end{proof}

Définissons, à partir de la forme de Smith et non de l'ordre analytique,
\begin{equation}\label{eq:pdf2-bsd-a-e-smith}
 a_E(T):=T^{-d}\det S_E(T).
\end{equation}
Par \eqref{eq:pdf2-bsd-rboc-lt}, $a_E$ est holomorphe dans $T=0$ et
\begin{equation}\label{eq:pdf2-bsd-a-e-cero}
 a_E(0)=R_{\rm Boc}^{\rm der}(S_E)\ne0.
\end{equation}
En substituant cette identité et
\eqref{eq:pdf2-bsd-t-e-construido} dans la factorisation analytique, on obtient
matériellement
\begin{equation}\label{eq:pdf2-bsd-expansion-local-deducida}
 \begin{aligned}
 L(E,s)
  &=(s-1)^d\,U_E^{\rm lead}(s),\\
 U_E^{\rm lead}(s)
  &:=u_E(s)
     \left(\frac{T_E(s)}{s-1}\right)^d
     a_E(T_E(s)),\\
 U_E^{\rm lead}(1)
  &=1\cdot1^d\cdot R_{\rm Boc}^{\rm der}(S_E)
   =R_{\rm Boc}^{\rm der}(S_E)\ne0.
 \end{aligned}
\end{equation}
L'exposant et le coefficient principal se déduisent donc du germe
holomorphe non nul ; ils ne sont pas des entrées de $T_E$, $u_E$ ou $S_E$.
Définissons maintenant, après cette déduction,
\begin{equation}\label{eq:pdf2-bsd-dan-zan}
 d_{\rm an}:=\operatorname{ord}_{s=1}L(E,s),
 \qquad
 \mathfrak z_{\rm an}:=
 \left[(s-1)^{-d_{\rm an}}L(E,s)\right]_{s=1}.
\end{equation}
L'équation \eqref{eq:pdf2-bsd-expansion-local-deducida} prouve
\begin{equation}\label{eq:pdf2-bsd-dan-d}
 d_{\rm an}=d,
 \qquad
 \mathfrak z_{\rm an}
 =R_{\rm Boc}^{\rm der}(S_E)
 \quad\text{sur la ligne basée}.
\end{equation}

Du côté arithmétique, le théorème~\ref{thm:pdf2-bsd-psi-pt} et la suite
\eqref{eq:pdf2-bsd-sha-tamagawa} construisent
\begin{equation}\label{eq:pdf2-bsd-dar-zar}
 \begin{aligned}
 d_{\rm ar}&:=\dim_K M_E^{\rm MW}=d,\\
 \mathfrak z_{\rm ar}&:=
 \frac{|\operatorname{Sha}(E)|\operatorname{Reg}(E)\prod_vc_v}
 {|E(\mathbb Q)_{\rm tors}|^2}
 \quad\text{sur la ligne arithmétique orientée}.
 \end{aligned}
\end{equation}

Pour typer les quatre flèches, si $C^\bullet$ est un complexe parfait à
base ordonnée $\mathcal B_i=(c_{i,1},\ldots,c_{i,r_i})$, on utilise
\begin{equation}\label{eq:pdf2-bsd-linea-determinante-base}
 \det C^\bullet:=
 \bigotimes_i\left(\bigwedge^{r_i}C^i\right)^{(-1)^i},
 \qquad
 \mathbf e_C:=\bigotimes_i
 (c_{i,1}\wedge\cdots\wedge c_{i,r_i})^{(-1)^i}.
\end{equation}
Pour un groupe fini $A$ avec une présentation basée
$\mathbb Z^t\xrightarrow{D_A}\mathbb Z^t\to A$, sa ligne
$\det_{\mathbb Z}A$ porte le générateur $\mathbf e_A$ et l'évaluation
\begin{equation}\label{eq:pdf2-bsd-evaluacion-grupo-finito}
 \operatorname{ev}_A(\mathbf e_A):=|\det D_A|=|A|.
\end{equation}
Lorsque cette ligne se combine à une ligne sur $K=\mathbb Q$, on utilise toujours
l'extension des scalaires
\begin{equation}\label{eq:pdf2-bsd-linea-finita-k}
 \det_KA:=\det_{\mathbb Z}A\otimes_{\mathbb Z}K;
\end{equation}
on ne tensorise pas directement des modules sur des anneaux distincts.
En particulier, on choisit les bases de Smith déjà construites dans
\eqref{eq:pdf2-bsd-smith-finito}, $m_1,\ldots,m_t$ de $M_E$ et
$n_1,\ldots,n_t$ de $N_E$, avec $D_Em_i=s_in_i$, et on fixe
\begin{equation}\label{eq:pdf2-bsd-generador-smith-fe}
 \mathbf e_{F_E}:=
 (m_1\wedge\cdots\wedge m_t)\otimes
 (n_1\wedge\cdots\wedge n_t)^{-1},
 \qquad
 \operatorname{ev}_{F_E}(\mathbf e_{F_E})=\prod_i|s_i|.
\end{equation}
Les lignes successives sont
\begin{equation}\label{eq:pdf2-bsd-cinco-lineas-tipadas}
\begin{aligned}
 \Delta_K&:=\det\mathcal C_E^K,\\
 \Delta_{\rm FERS}&:=\det\mathcal C_E^{\rm Iw},\\
 \Delta_{\rm PT}&:=
   \det_K M_E^{\rm MW}\otimes
   \det_K(M_E^{\rm MW})^\vee\otimes\det_KF_E,\\
 \Delta_{\rm STS}&:=
   \det_K M_E^{\rm MW}\otimes
   \det_K(M_E^{\rm MW})^\vee\otimes
   \det_K\operatorname{Sha}(E)
   \otimes\bigotimes_{v<\infty}
      \det_K\Phi_v(\mathbb F_v),\\
 \Delta_{\rm BSD}&:=\mathbb R\,\mathbf e_{\rm BSD}.
\end{aligned}
\end{equation}
La spécialisation principale n'est pas une identification tacite. Notons
$I:=(T)\subset K[[T]]$ et
$\mathfrak c_{K/{\rm FERS}}^{\rm reg}:=
\det_{\rm alt}\Phi_{K/{\rm FERS}}:\Delta_K[[T]]\to
\Delta_{\rm FERS}[[T]]$. Soient
\begin{equation}\label{eq:pdf2-bsd-lineas-leading}
 \begin{aligned}
  \Delta_{\rm FERS}^{(d)}&:=I^d\Delta_{\rm FERS}[[T]],&
  \Delta_{\rm FERS}^{\rm lead}
    &:=I^{-d}\Delta_{\rm FERS}^{(d)}
       \otimes_{K[[T]],\,T\mapsto0}K,\\
  \Delta_K^{(d)}&:=(\mathfrak c_{K/{\rm FERS}}^{\rm reg})^{-1}
       (\Delta_{\rm FERS}^{(d)}),&
  \Delta_K^{\rm lead}
    &:=I^{-d}\Delta_K^{(d)}
       \otimes_{K[[T]],\,T\mapsto0}K.
 \end{aligned}
\end{equation}
Ici, $I^{-d}\Delta^{(d)}$ désigne le module formé par les symboles
$T^{-d}z$ avec $z\in\Delta^{(d)}$. La spécialisation n'est définie
que sur la partie d'ordre au moins $d$ :
\begin{equation}\label{eq:pdf2-bsd-sp-leading-tipado}
 \operatorname{sp}^{(d)}_{\rm FERS}:
 \Delta_{\rm FERS}^{(d)}\longrightarrow
 \Delta_{\rm FERS}^{\rm lead},
 \qquad
 z(T)\longmapsto
 \overline{T^{-d}z(T)}.
\end{equation}
L'équivalence Kato--FERS induit alors l'application typée
\begin{equation}\label{eq:pdf2-bsd-c-k-fers-leading}
 \begin{aligned}
 \mathfrak c_{K/{\rm FERS}}^{\rm lead}:{}
 &\Delta_K^{\rm lead}\xrightarrow{\ \sim\ }
   \Delta_{\rm FERS}^{\rm lead},\\
 &\overline{T^{-d}z_K(T)}\longmapsto
   \overline{T^{-d}
   \mathfrak c_{K/{\rm FERS}}^{\rm reg}(z_K(T))}.
 \end{aligned}
\end{equation}
La formule n'applique pas $\operatorname{sp}^{(d)}$ à une ligne nue. En
particulier, si
\begin{equation}\label{eq:pdf2-bsd-seccion-zeta-leading}
 \begin{aligned}
  \boldsymbol\zeta_{\rm FERS}(T)
    &:=\det S_E(T)\,\mathbf e_{\rm FERS}
       \in\Delta_{\rm FERS}^{(d)},\\
  \boldsymbol\zeta_K(T)
    &:=(\mathfrak c_{K/{\rm FERS}}^{\rm reg})^{-1}
       \boldsymbol\zeta_{\rm FERS}(T)
       \in\Delta_K^{(d)},
 \end{aligned}
\end{equation}
alors
\begin{equation}\label{eq:pdf2-bsd-especializacion-zeta}
 \operatorname{sp}^{(d)}_{\rm FERS}
       (\boldsymbol\zeta_{\rm FERS})
 =R_{\rm Boc}^{\rm der}(S_E)\,
       \mathbf e_{\rm FERS}^{\rm lead},
 \quad
 \mathfrak c_{K/{\rm FERS}}^{\rm lead}
       (\overline{T^{-d}\boldsymbol\zeta_K})
 =R_{\rm Boc}^{\rm der}(S_E)\,
       \mathbf e_{\rm FERS}^{\rm lead}.
\end{equation}
La définition utilise uniquement $d=\operatorname{ord}_T\det S_E(T)$, déjà calculé par
Smith, et non l'ordre analytique.

Les quatre flèches propriétaires, maintenant dotées de domaines et de codomaines explicites,
sont comparées dans la fibre réelle dès le premier endroit où intervient la
hauteur de Néron. Pour éviter de multiplier une ligne sur $K$ par un
scalaire réel n'appartenant pas en général à $K$, on fixe
\begin{equation}\label{eq:pdf2-bsd-lineas-reales-comparador}
 \begin{aligned}
  \Delta_{K,\mathbb R}^{\rm lead}
    &:=\Delta_K^{\rm lead}\otimes_{K,\iota_\infty}\mathbb R,&
  \Delta_{{\rm FERS},\mathbb R}^{\rm lead}
    &:=\Delta_{\rm FERS}^{\rm lead}\otimes_{K,\iota_\infty}\mathbb R,\\
  \Delta_{{\rm PT},\mathbb R}
    &:=\Delta_{\rm PT}\otimes_{K,\iota_\infty}\mathbb R,&
  \Delta_{{\rm STS},\mathbb R}
    &:=\Delta_{\rm STS}\otimes_{K,\iota_\infty}\mathbb R.
 \end{aligned}
\end{equation}
Les quatre flèches sont alors
\begin{equation}\label{eq:pdf2-bsd-cuatro-comparadores}
\begin{aligned}
 \mathfrak c_{K/{\rm FERS},\mathbb R}^{\rm lead}:
   &\Delta_{K,\mathbb R}^{\rm lead}
      \longrightarrow\Delta_{{\rm FERS},\mathbb R}^{\rm lead},
 &\mathfrak c_{K/{\rm FERS},\mathbb R}^{\rm lead}
   &:=\mathfrak c_{K/{\rm FERS}}^{\rm lead}\otimes_K1_{\mathbb R},\\
 \mathfrak c_{{\rm PT},\mathbb R}:
   &\Delta_{{\rm FERS},\mathbb R}^{\rm lead}
      \longrightarrow\Delta_{{\rm PT},\mathbb R},
 &\mathfrak c_{{\rm PT},\mathbb R}(z)&:=
   \mathfrak n_{E,\mathbb R}(z)\otimes\mathbf e_{F_E},\\
 \mathfrak c_{{\rm STS},\mathbb R}:
   &\Delta_{{\rm PT},\mathbb R}
      \longrightarrow\Delta_{{\rm STS},\mathbb R},
 &\mathfrak c_{{\rm STS},\mathbb R}&:=\det\bigl(
   0\to\operatorname{Sha}(E)\to F_E
     \to\textstyle\bigoplus_{v<\infty}\Phi_v(\mathbb F_v)
     \to0\bigr)\otimes_K1_{\mathbb R},\\
 \mathfrak c_{\rm tors,\infty}:
   &\Delta_{{\rm STS},\mathbb R}
      \longrightarrow\Delta_{\rm BSD}.&&
\end{aligned}
\end{equation}

\begin{proposition}[Calcul générateur par générateur des quatre flèches]
\label{prop:pdf2-bsd-cuatro-flechas-generadores}
Les actions sur les bases de
\eqref{eq:pdf2-bsd-cinco-lineas-tipadas} sont les suivantes.

\emph{Première flèche.} Si
$\mathbf e_K^i=\bigwedge_{b\in\mathcal B_i^K}b$ et
$\mathbf e_{\rm FERS}^i=\bigwedge_{b\in\mathcal B_i^K}\upsilon_i(b)$,
la table \eqref{eq:pdf2-bsd-tabla-kato-fers} donne
\begin{equation}\label{eq:pdf2-bsd-c1-generadores}
 \mathfrak c_{K/{\rm FERS}}^{\rm reg}(\mathbf e_K)
 =\prod_i\det(I+TB_i(T))^{(-1)^i}\mathbf e_{\rm FERS}
 =\rho_E(T)\mathbf e_{\rm FERS},
 \qquad \rho_E(0)=1.
\end{equation}
Si $\mathbf e_K^{\rm lead}$ et $\mathbf e_{\rm FERS}^{\rm lead}$ sont les
bases induites par \eqref{eq:pdf2-bsd-lineas-leading}, alors
\begin{equation}\label{eq:pdf2-bsd-c1-leading-generador}
 \mathfrak c_{K/{\rm FERS},\mathbb R}^{\rm lead}
     (\mathbf e_K^{\rm lead}\otimes1)
 =\rho_E(0)(\mathbf e_{\rm FERS}^{\rm lead}\otimes1)
 =\mathbf e_{\rm FERS}^{\rm lead}\otimes1.
\end{equation}

\emph{Deuxième flèche.} Choisissons des bases intégrales
$a_{q,1},\ldots,a_{q,r_q}$ de $A_{q,\mathbb Z}$, vues dans $A_q$, et
$b_{q,j}=\bar\beta_q(a_{q,j})$ de $B_q$. La matrice de la hauteur de la page
est définie par
\begin{equation}\label{eq:pdf2-bsd-hq-generadores}
 j_q(b_{q,j})=
 \sum_{k=1}^{r_q}(H_q)_{kj}\,a_{q,k}^\vee\otimes T^q.
\end{equation}
Par conséquent, sur le produit extérieur complet,
\begin{equation}\label{eq:pdf2-bsd-c2-cuna-paginas}
 \bigwedge_{j=1}^{r_q}j_q\bar\beta_q(a_{q,j})
 =T^{qr_q}\det(H_q)
  \bigwedge_{k=1}^{r_q}a_{q,k}^\vee.
\end{equation}
Le produit sur $q$, avec les lignes régulières, a pour exposant
$\sum_q qr_q=\sum_q\dim_K V_q=d$. Pour
$1\leq j\leq r_q$ et $1\leq k\leq q$, posons
$e_{q,j,k}:=a_{q,j}e_{q,k}$ dans
$\operatorname{Jet}_{\rm Boc,\mathbb Z}(E)$ et définissons
\begin{equation}\label{eq:pdf2-bsd-base-mw-desde-jets}
 P_{q,j,k}:=\Psi_{{\rm PT},\mathbb Z}^{-1}
       \mathcal U_{\rm jet}(e_{q,j,k})
 \in M_{E,\mathbb Z}^{\rm MW}.
\end{equation}
Par \eqref{eq:pdf2-bsd-desdoblamiento-jets},
\eqref{eq:pdf2-bsd-dimension-jet} et
\eqref{eq:pdf2-bsd-psi-pt-integral}, ces $d$ éléments sont une base
intégrale, non un sous-réseau d'indice non contrôlé. Ils sont énumérés dans l'ordre
lexicographique comme $P_1,\ldots,P_d$ et, sur $K$, forment la base de
$M_E^{\rm MW}$ utilisée ensuite.

La hauteur est construite indépendamment de FERS.\@ Pour
$P\in E(\overline{\mathbb Q})$, on fixe
\begin{equation}\label{eq:pdf2-bsd-altura-canonica-construida}
 \widehat h_E(P):=
 \frac12\lim_{n\to\infty}4^{-n}h_x([2^n]P),
 \qquad
 \langle P,Q\rangle_{\rm NT}:=
 \frac12\bigl(\widehat h_E(P+Q)-\widehat h_E(P)-\widehat h_E(Q)\bigr).
\end{equation}
De manière équivalente, le biextenseur de Poincaré rigidifié sur le modèle de
Néron produit, pour toute place $v$, y compris $v=\infty$, le symbole local
réel $\lambda_v(P,Q)$ normalisé par la section nulle, et
\begin{equation}\label{eq:pdf2-bsd-altura-neron-local-global}
 \langle P,Q\rangle_{\rm NT}
 =\sum_v\lambda_v(P,Q)\in\mathbb R.
\end{equation}
Cette hauteur réelle, symétrique et positive n'est pas remplacée par une somme
d'invariants de cup-produits à valeurs dans $\mathbb Q/\mathbb Z$.

Posons
$H_{\rm NT}^{\flat}(P)(Q):=\langle P,Q\rangle_{\rm NT}$ et soit
$\mathcal B_E:=\Psi_{{\rm PT},\mathbb Z}^{-1}\mathcal U_{\rm jet}$,
étendu à $\mathbb R$. Les $q-1$ directions supplémentaires ne sont pas remplacées
par des identités : on calcule le bloc de Gram complet, y compris tous les
croisements entre profondeurs, lignes et jets,
\begin{equation}\label{eq:pdf2-bsd-gram-jet-completo}
 \begin{aligned}
 G_E^{\rm jet}\bigl((q,j,k),(q',j',k')\bigr)
 &:={}
 \left\langle P_{q,j,k},P_{q',j',k'}\right\rangle_{\rm NT}\\
 &=\sum_v\lambda_v
   \left(P_{q,j,k},P_{q',j',k'}\right),
 \end{aligned}
\end{equation}
pour tout $1\le k\le q$, $1\le k'\le q'$. Cette matrice est matérielle : les
$P_{q,j,k}$ ont été construits dans
\eqref{eq:pdf2-bsd-base-mw-desde-jets} et chaque entrée est calculée par
\eqref{eq:pdf2-bsd-altura-neron-local-global}. Par définition du pullback
d'une forme bilinéaire,
\begin{equation}\label{eq:pdf2-bsd-altura-nt-generadores}
 \boxed{
 \mathcal B_E^\vee H_{\rm NT}^{\flat}\mathcal B_E=G_E^{\rm jet}.}
\end{equation}
On n'affirme pas que $G_E^{\rm jet}$ est diagonale par blocs. Dans la base
intégrale $P_1,\ldots,P_d$, le régulateur et son action sur la ligne déterminante
sont exactement
\begin{equation}\label{eq:pdf2-bsd-regulador-cuna-altura}
 \begin{aligned}
 \operatorname{Reg}(E)&:=\det G_E^{\rm jet}>0,\\
 \bigwedge_{j=1}^dH_{\rm NT}^{\flat}(P_j)
 &=\operatorname{Reg}(E)
   (P_1^\vee\wedge\cdots\wedge P_d^\vee).
 \end{aligned}
\end{equation}
La page de Bockstein et la hauteur sont deux accouplements construits, avec des
bases déjà fixées. Le comparateur de Néron entre leurs lignes est défini avant
d'utiliser le terme principal par
\begin{equation}\label{eq:pdf2-bsd-comparador-neron-jet}
 \begin{aligned}
 \mathfrak n_{E,\mathbb R}:
 \Delta_{{\rm FERS},\mathbb R}^{\rm lead}&\longrightarrow
 \bigl(\det_K(M_E^{\rm MW})\otimes_K
 \det_K(M_E^{\rm MW})^\vee\bigr)\otimes_{K,\iota_\infty}\mathbb R,\\
 \mathfrak n_{E,\mathbb R}
 (\mathbf e_{\rm FERS}^{\rm lead}\otimes1)
 &:=\frac{\operatorname{Reg}(E)}{R_{\rm Boc}^{\rm der}(S_E)}
 (P_1\wedge\cdots\wedge P_d)
 \otimes(P_1^\vee\wedge\cdots\wedge P_d^\vee)\otimes1.
 \end{aligned}
\end{equation}
Le dénominateur est non nul par \eqref{eq:pdf2-bsd-a-e-cero} ; numérateur et
dénominateur s'obtiennent, respectivement, à partir du Gram complet
\eqref{eq:pdf2-bsd-gram-jet-completo} et des pages
\eqref{eq:pdf2-bsd-rboc-lt}. Ainsi, la normalisation est un quotient réel de deux
déterminants préalablement calculés, non une orthogonalité inventée ni une
consultation de la formule BSD. L'extension des scalaires s'effectue avant
d'appliquer le régulateur et se conserve dans les flèches PT et STS.

Les pages supérieures et la ligne régulière conservent le générateur intégral de
$F_E$ au moyen d'une application typée. On utilise l'extension rationnelle de la section
$s_q$ déjà construite dans \eqref{eq:pdf2-bsd-desdoblamiento-jets}. La section duale
\begin{equation}\label{eq:pdf2-bsd-glue-dual-tipado}
 \operatorname{gl}_q^\vee:A_q^\vee\longrightarrow
   (\operatorname{Surv}_{\rm det}^{\rm MW})^\vee
\end{equation}
est définie dans la base complète par
$(\operatorname{gl}_q^\vee\varphi)
(\operatorname{gl}_r s_r(a))=\delta_{qr}\varphi(a)$ et par zéro dans les lignes
FS.\@ Soit $\mathcal I_{\rm fin}$ la liste ordonnée des lignes régulières et des paires
$(q,j)$ avec $q\geq2$. Pour $\eta$ avec $q(\eta)\geq2$, on pose
\begin{equation}\label{eq:pdf2-bsd-lifts-filas-smith}
 \begin{aligned}
 \widetilde a_\eta&:=\operatorname{pr}_{M_E}\operatorname{pr}_{\mathcal L_E}
   \Theta_E\operatorname{gl}_{q(\eta)}s_{q(\eta)}(a_\eta),\\
 \widetilde b_\eta^\vee&:=
   \operatorname{pr}_{N_E^\vee}\operatorname{pr}_{\mathcal L_E^\vee}
   (\Theta_E^\vee)^{-1}\operatorname{gl}_{q(\eta)}^\vee
   \bigl((1\otimes\partial_T^{(q(\eta))})
      j_{q(\eta)}\bar\beta_{q(\eta)}(a_\eta)\bigr).
 \end{aligned}
\end{equation}
Dans une ligne régulière, on définit séparément
\begin{equation}\label{eq:pdf2-bsd-lifts-fila-regular}
 \widetilde a_\eta:=
 \operatorname{pr}_{M_E}\operatorname{pr}_{\mathcal L_E}
 \Theta_Es_{\rm reg}(a_\eta),
 \qquad
 \widetilde b_\eta^\vee:=
 \operatorname{pr}_{N_E^\vee}\operatorname{pr}_{\mathcal L_E^\vee}
 (\Theta_E^\vee)^{-1}s_{\rm reg}^\vee
 (\bar S_E(0)a_\eta).
\end{equation}
Tous les domaines sont maintenant $A_q$ ou $A_q^\vee$ selon le cas ; on
n'applique pas $\operatorname{gl}_q$ à un covecteur.

La bijection requise ne se déduit pas de matrices unimodulaires. Les
équations d'unité--counité montrent sur chaque générateur que les listes
$(\widetilde a_\eta)_\eta$ et leurs terminaux duaux sont des bases de $M_E$ et
$N_E^\vee$. Écrivons la liste héritée comme
$\mathcal I_{\rm fin}=\{\eta_1<\cdots<\eta_t\}$ et définissons directement
\begin{equation}\label{eq:pdf2-bsd-nu-construida}
 \boxed{\nu(\eta_j):=j},
 \qquad
 \nu:\mathcal I_{\rm fin}\xrightarrow{\sim}\{1,\ldots,t\}.
\end{equation}
Posons $\mu_j:=\widetilde a_{\eta_j}$ et soit
$\xi_1,\ldots,\xi_t$ la base terminale de $N_E$. Avant Smith, la matrice
matérielle est
\begin{equation}\label{eq:pdf2-bsd-c2-filas-finitas}
 D^{\rm term}_{kj}:=
 \langle\xi_k^\vee,D_E\mu_j\rangle,
 \qquad
 \widetilde b_{\eta_j}^\vee
 =\sum_{k=1}^tD^{\rm term}_{kj}\xi_k^\vee.
\end{equation}
Cette égalité est la relation terminale de
\eqref{eq:pdf2-bsd-descomposicion-global}, évaluée dans la counité
$\xi_k^\vee$ ; elle ne divise pas par un diviseur élémentaire. Si
$UD^{\rm term}V=\operatorname{diag}(s_1,\ldots,s_t)$, alors
\begin{equation}\label{eq:pdf2-bsd-cuna-smith-con-uv}
 \bigwedge_{j=1}^t\widetilde b_{\eta_j}^\vee
 =\det(D^{\rm term})
   (\xi_1^\vee\wedge\cdots\wedge\xi_t^\vee),
 \qquad
 |\det(D^{\rm term})|=\prod_{j=1}^t|s_j|.
\end{equation}
Les facteurs $\det U,\det V\in\{\pm1\}$ sont enregistrés par
$\operatorname{or}_{\rm det}$ lorsqu'on passe de $(\mu_j),(\xi_j)$ aux bases de
Smith $(m_j),(n_j)$. Le quotient des produits extérieurs est donc exactement
$\mathbf e_{F_E}$ de \eqref{eq:pdf2-bsd-generador-smith-fe} ; on n'a pas
supposé que $U$ ou $V$ soient des matrices de permutation et aucune page
ne reste non publiée. En
conséquence,
\begin{equation}\label{eq:pdf2-bsd-c2-generadores}
 \mathfrak c_{{\rm PT},\mathbb R}
 \bigl(R_{\rm Boc}^{\rm der}(S_E)
       (\mathbf e_{\rm FERS}^{\rm lead}\otimes1)\bigr)
 =\operatorname{Reg}(E)\,
  (P_1\wedge\cdots\wedge P_d)\otimes
  (P_1^\vee\wedge\cdots\wedge P_d^\vee)
  \otimes\mathbf e_{F_E}\otimes1.
\end{equation}

\emph{Troisième flèche.} Pour chaque $\ell$, choisissons des générateurs ordonnés
$\sigma_{\ell,a}$ de $\operatorname{Sha}(E)[\ell^\infty]$ et
$\phi_{\ell,v,b}$ de $\Phi_v(\mathbb F_v)[\ell^\infty]$. Les relèvements explicites
$\operatorname{Lift}_{\Phi,\ell}(\phi_{\ell,v,b})$ et les images
$\iota_{{\rm Sha},\ell}(\sigma_{\ell,a})$ forment, par exactitude, une base de
présentation de $F_{E,\ell}$. Dans ces bases, la matrice des relations se
calcule générateur par générateur comme
\begin{equation}\label{eq:pdf2-bsd-c3-matriz-presentacion}
 D_{F,\ell}=
 \begin{pmatrix}
  D_{{\rm Sha},\ell}&X_\ell\\
  0&\displaystyle\bigoplus_{v<\infty}D_{\Phi_v,\ell}
 \end{pmatrix}.
\end{equation}
La colonne $X_\ell$ enregistre le changement de chaque relèvement lorsqu'on modifie son
représentant ; le bloc inférieur gauche est nul car
$\partial_{{\rm loc},\ell}\iota_{{\rm Sha},\ell}=0$. Le produit extérieur de
présentation ne dépend donc pas de $X_\ell$ et
$\det D_{F,\ell}=\det D_{{\rm Sha},\ell}
\prod_{v<\infty}\det D_{\Phi_v,\ell}$. Sur celui-ci,
\begin{equation}\label{eq:pdf2-bsd-c3-generadores}
 \mathbf e_{F_{E,\ell}}
 \longmapsto
 \mathbf e_{\operatorname{Sha}(E)[\ell^\infty]}
 \otimes\bigotimes_{v<\infty}
    \mathbf e_{\Phi_v(\mathbb F_v)[\ell^\infty]},
\end{equation}
et les évaluations de \eqref{eq:pdf2-bsd-evaluacion-grupo-finito} donnent
$|F_{E,\ell}|=|\operatorname{Sha}(E)[\ell^\infty]|
\prod_{v<\infty}|\Phi_v(\mathbb F_v)[\ell^\infty]|$.
Multiplier sur $\ell$ produit
\begin{equation}\label{eq:pdf2-bsd-c3-integral-generadores}
 \mathfrak c_{{\rm STS},\mathbb R}(\mathbf e_{F_E}\otimes1)
 =(\mathbf e_{\operatorname{Sha}(E)}
  \otimes\bigotimes_{v<\infty}\mathbf e_{\Phi_v(\mathbb F_v)})\otimes1,
 \qquad
 |F_E|=|\operatorname{Sha}(E)|\prod_vc_v.
\end{equation}

\emph{Quatrième flèche.} La base de Mordell--Weil précédente et une
différentielle de Néron $\omega_E$ de période positive
$\Omega_E$ étant fixées, la base orientée de $\Delta_{\rm STS}$ est envoyée par
\begin{equation}\label{eq:pdf2-bsd-c4-generadores}
 \begin{aligned}
 &x\,(P_1\wedge\cdots\wedge P_d)\otimes
 (P_1^\vee\wedge\cdots\wedge P_d^\vee)\otimes
 \mathbf e_{\operatorname{Sha}(E)}
 \otimes\bigotimes_{v<\infty}\mathbf e_{\Phi_v(\mathbb F_v)}\\
 &\hspace{28mm}\longmapsto
 \frac{x\,\Omega_E\,
   \operatorname{ev}_{\operatorname{Sha}(E)}
      (\mathbf e_{\operatorname{Sha}(E)})
   \prod_{v<\infty}\operatorname{ev}_{\Phi_v(\mathbb F_v)}
      (\mathbf e_{\Phi_v(\mathbb F_v)})}
 {|E(\mathbb Q)_{\rm tors}|^2}\,
 \mathbf e_{\rm BSD}.
\end{aligned}
\end{equation}
Par \eqref{eq:pdf2-bsd-evaluacion-grupo-finito}, le coefficient de la dernière
ligne est
$x\Omega_E|\operatorname{Sha}(E)|\prod_vc_v/
|E(\mathbb Q)_{\rm tors}|^2$.
Les deux copies du groupe de torsion proviennent des degrés duaux $0$ et $2$ ;
le facteur $\Omega_E$ est l'évaluation de $\omega_E$ sur la composante réelle
orientée. Cette orientation place $\Omega_E\mathfrak z_{\rm ar}$, et non son
inverse, comme coefficient arithmétique de la composition ; la comparaison avec
$\mathfrak z_{\rm an}$ ne sera déduite qu'ensuite au moyen du carré des
trivialisations.
\end{proposition}

\begin{proof}
La première identité est le déterminant alterné des matrices triangulaires
$I+TB_i(T)$. La deuxième applique la spécialisation des pages puis
le comparateur de Néron \eqref{eq:pdf2-bsd-comparador-neron-jet}. Son
numérateur est le déterminant du Gram de jets complet
\eqref{eq:pdf2-bsd-gram-jet-completo}, avec tous les croisements, et son dénominateur
est le déterminant de Bockstein déjà calculé. Par
\eqref{eq:pdf2-bsd-regulador-cuna-altura}, son action sur les produits extérieurs est
\eqref{eq:pdf2-bsd-c2-generadores} ; l'exposant $d$ s'obtient avant de
prendre les dimensions de Mordell--Weil. La troisième est l'isomorphisme canonique de lignes associé à la
suite exacte, écrit dans les sections
$\iota_{{\rm Sha},\ell}$ et $\operatorname{Lift}_{\Phi,\ell}$ déjà
construites. La quatrième est l'évaluation des deux complexes finis de
torsion et de la ligne de Néron. Cela calcule les quatre flèches sans utiliser
l'égalité BSD.
\end{proof}

Leur composition
\begin{equation}\label{eq:pdf2-bsd-comparador-total}
 \mathfrak C_E:=
 \mathfrak c_{\rm tors,\infty}\circ
 \mathfrak c_{{\rm STS},\mathbb R}\circ
 \mathfrak c_{{\rm PT},\mathbb R}\circ
 \mathfrak c_{K/{\rm FERS},\mathbb R}^{\rm lead}:
 \Delta_{K,\mathbb R}^{\rm lead}
 \longrightarrow\Delta_{\rm BSD}
\end{equation}
est bien typée. Avant d'insérer l'élément analytique ou quelque cardinal que ce soit,
on fixe les trivialisations indépendantes
\begin{equation}\label{eq:pdf2-bsd-trivializaciones-independientes}
 \operatorname{tr}_{\rm an}(x\mathbf e_K^{\rm lead}):=x,
 \qquad
 \operatorname{tr}_{\rm Ner}(y\mathbf e_{\rm BSD}):=y.
\end{equation}
La première utilise uniquement la base de Kato spécialisée par
\eqref{eq:pdf2-bsd-lineas-leading} et la normalisation Mellin--Fredholm de
\eqref{eq:pdf2-bsd-fredholm-euler} ; en particulier,
$\operatorname{tr}_{\rm an}
(\overline{T^{-d}\boldsymbol\zeta_K})
=R_{\rm Boc}^{\rm der}(S_E)=\mathfrak z_{\rm an}$ par
\eqref{eq:pdf2-bsd-especializacion-zeta}. La deuxième utilise uniquement la différentielle
de Néron orientée et satisfait
$\operatorname{tr}_{\rm Ner}(\mathbf e_{\rm BSD})=1$. L'identité du
comparateur basé de $\mathfrak G_{\rm det}(E)$, déployée par les quatre
flèches, est le
carré
\begin{equation}\label{eq:pdf2-bsd-cuadrado-trivializaciones}
\begin{array}{ccc}
 \Delta_{K,\mathbb R}^{\rm lead}
   &\xrightarrow{\ \mathfrak C_E\ }&\Delta_{\rm BSD}\\[2mm]
 \big\downarrow\scriptstyle\operatorname{tr}_{\rm an}
   &&\big\downarrow\scriptstyle\operatorname{tr}_{\rm Ner}\\[2mm]
 \mathbb R&=&\mathbb R.
\end{array}
\end{equation}
Ce carré est fixé avec les bases structurelles avant d'évaluer
$L(E,s)$, le régulateur ou les groupes finis.

L'élément analytique est placé, après construction des flèches, dans son
domaine correct :
\begin{equation}\label{eq:pdf2-bsd-elemento-analitico-tipado}
 \mathbf z_E^{\rm an}:=\mathfrak z_{\rm an}\mathbf e_K^{\rm lead}
 =\overline{T^{-d}\boldsymbol\zeta_K}
 \in\Delta_{K,\mathbb R}^{\rm lead}.
\end{equation}
Par \eqref{eq:pdf2-bsd-dan-d},
$\mathfrak z_{\rm an}=R_{\rm Boc}^{\rm der}(S_E)$ ; la deuxième égalité de
\eqref{eq:pdf2-bsd-especializacion-zeta} prouve que la première flèche envoie
$\mathbf z_E^{\rm an}$ exactement sur le membre gauche de
\eqref{eq:pdf2-bsd-c2-generadores}. Les trois autres flèches, calculées sur les
générateurs, donnent
\begin{equation}\label{eq:pdf2-bsd-elementos-linea}
 d_{\rm an}=d=d_{\rm ar},
 \qquad
 \boxed{
 \mathfrak C_E(\mathbf z_E^{\rm an})
 =\Omega_E\mathfrak z_{\rm ar}\,\mathbf e_{\rm BSD}.}
\end{equation}
Appliquer les flèches verticales de
\eqref{eq:pdf2-bsd-cuadrado-trivializaciones} produit l'égalité numérique
\begin{equation}\label{eq:pdf2-bsd-igualdad-numerica-deducida}
 \boxed{\mathfrak z_{\rm an}=\Omega_E\mathfrak z_{\rm ar}.}
\end{equation}

\begin{theorem}[Publication de Birch--Swinnerton--Dyer]
\label{thm:pdf2-bsd-publicacion-final}
Pour l'image déterminantale autonome complète
$\mathfrak G_{\rm det}(E)$, on a
\begin{equation}\label{eq:pdf2-bsd-rango-final}
 \operatorname{ord}_{s=1}L(E,s)=\operatorname{rank}E(\mathbb Q)=:r
\end{equation}
et
\begin{equation}\label{eq:pdf2-bsd-formula-final}
 \frac{L^{(r)}(E,1)}{r!\,\Omega_E}
 =\frac{|\operatorname{Sha}(E)|\operatorname{Reg}(E)\prod_vc_v}
 {|E(\mathbb Q)_{\rm tors}|^2}.
\end{equation}
\end{theorem}

\begin{proof}
L'égalité des degrés est
\eqref{eq:pdf2-bsd-dan-d} suivie de
\eqref{eq:pdf2-bsd-rango-orden}. La table Kato--FERS transporte l'élément
analytique vers toutes les pages de Bockstein ; $\Psi_{\rm PT}$ les recolle au réseau
de Mordell--Weil et publie son déterminant de hauteur ; la suite exacte
Sha--Tamagawa décompose le facteur de Smith ; et le dernier comparateur insère
le réseau torsionnel et la période déjà orientée. C'est exactement
\eqref{eq:pdf2-bsd-elementos-linea}. Le carré indépendant
\eqref{eq:pdf2-bsd-cuadrado-trivializaciones} donne
\eqref{eq:pdf2-bsd-igualdad-numerica-deducida} ; substituer les définitions de
$\mathfrak z_{\rm an}$ et $\mathfrak z_{\rm ar}$ produit
\eqref{eq:pdf2-bsd-formula-final}. Aucune des quatre flèches n'est
normalisée avec $r$, $\operatorname{Reg}(E)$,
$|\operatorname{Sha}(E)|$ ou le terme principal : ces valeurs n'apparaissent
qu'à la fin de leurs déterminants respectifs.
\end{proof}

\section{Contrôles non directeurs, réfutateurs et statut}

\begin{falsifier}[Réfutateurs séparés par composante]
La construction échoue si une histoire de
$\Sigma_{\rm det}^{\rm mult}$ n'appartient pas à la partition
\eqref{eq:pdf2-bsd-tres-proyectores} ; si l'on écrit
$\partial H_{\rm FS}+H_{\rm FS}\partial=q_{\rm root}$ et que l'on efface ainsi
$\mathcal L_E$ ; si une ligne non racine de
\eqref{eq:pdf2-bsd-tabla-kato-fers} a un poids nul ; si les équations
d'unité--counité \eqref{eq:pdf2-bsd-unidad-counidad-paginas} échouent ; si
$D_E\otimes\mathbb Q$ n'est pas inversible ; si
$\pi_\Phi\operatorname{Lift}_\Phi\ne I$ ; ou si une flèche de
\eqref{eq:pdf2-bsd-cuatro-comparadores} utilise la valeur finale pour fixer sa
base. Chaque réfutateur rompt une arête concrète et ne rouvre pas les précédentes.
\end{falsifier}

L'homotopie $H_{\rm FS}$ est un contrôle ultérieur du remplisseur explicite
$h_*=-vf$ et non une prémisse directrice. Les formes de Smith, les
déterminants et les cardinaux sont calculés après la construction des applications ;
ils ne remplacent pas les tables, le recollement des pages ni l'exactitude des cocycles.

\paragraph{Statut et provenance.}
Alexander--Whitney, le produit fibré, l'unité commune, la dualité
réduite, la rigidité de la racine, le remplisseur, toutes les pages de Bockstein et
les quatre lecteurs appartiennent à
\texttt{ARCHITECTURE\_AUTORALE\_PRÉEXISTANTE}/\texttt{EXACT\_INTERNE} de
$\mathfrak G_{\rm det}(E)$. La partition exhaustive
\eqref{eq:pdf2-bsd-descomposicion-global}, la table
\eqref{eq:pdf2-bsd-tabla-kato-fers}, les formules
\eqref{eq:pdf2-bsd-psi-y-lambda-q} et les applications de cocycles
\eqref{eq:pdf2-bsd-iota-sha-cociclos}--
\eqref{eq:pdf2-bsd-lift-componentes} sont
\texttt{FORMALISATION\_NOUVELLE} : elles déploient matériellement les entrelaceurs
déjà contenus dans les treize coordonnées structurelles. Le théorème final a la
force
\texttt{DÉMONTRÉ\_AVEC\_STRUCTURE\_DE\_DÉPART\_EXPLICITE} ; il ne reste pas
d'obligation résiduelle directrice et on n'utilise pas d'homotopie abstraite pour
fabriquer le résultat.
